\chapter{Fähigkeiten und Klassen}

Dieses Kapitel beschreibt alle spielrelevanten \term{Fähigkeiten}{faehigkeiten} eines Charakters:
Skills und Zauber, ihre Steigerung, Kosten, Einsatz im Spiel sowie ihre
legendären Erweiterungen.

\section{Grundlagen der Fähigkeiten}

Fähigkeiten sind aktive Spieloptionen, mit denen ein Charakter gezielt
Handlungen im Kampf, in Erkundung oder sozialen Situationen ausführt.
Es wird zwischen \textbf{Skills} (nichtmagisch) und \textbf{Zaubern} (magisch) unterschieden.

\subsection*{Skills und Zauber}

\begin{itemize}
	\item \term{Skills}{skills} sind körperliche, taktische oder handwerkliche Fähigkeiten
	und verbrauchen in der Regel \textbf{\termref[Aktionspunkte (AP)]{aktionspunkte-ap}}.
	\item \term{Zauber}{zauber} sind magische Fähigkeiten und verbrauchen
	\textbf{\termref[Manapunkte (MP)]{manapunkte-mp}}.
\end{itemize}

Das Ausführen einer Fähigkeit kann neben Aktionspunkte (AP) und Manapunkte (MP) auch zusätzliche Ressourcen (z.\,B. Materialien oder Lebenspunkte (LP)) verbrauchen. Hat ein Charakter \textbf{nicht genügend \termref[Ressourcen]{ressourcen}, kann die Fähigkeit nicht eingesetzt werden}.

Welche Art von Fähigkeiten ein Charakter lernen kann, hängt von seiner Klasse ab:
\begin{itemize}
	\item \termref[Nichtmagische Klassen]{nicht-magische-klassen} lernen Skills
	\item \termref[Magierklassen]{magierklassen} lernen Zauber
	\item Basis-Fähigkeiten können von allen Klassen erlernt werden
\end{itemize}

\subsection*{Fähigkeitspunkte}

Fähigkeiten werden mit \term{Fähigkeitspunkten}{faehigkeitspunkte}
(Skillpunkte oder Zauberpunkte) erlernt und gesteigert.

\begin{itemize}
	\item Eine neue Fähigkeit wird mit \textbf{1 Punkt} auf Rang 1 erlernt.
	\item Jeder weitere Rang kostet \textbf{1 Fähigkeitspunkt pro Rang}.
	\item Der maximale reguläre Rang beträgt \textbf{Rang 5}.
\end{itemize}

Fähigkeitspunkte sind \textbf{keine Verbrauchsressource im Spiel},
sondern dienen ausschließlich der Charakterentwicklung.

\subsection*{Rang und Skalierung}

Der Rang einer Fähigkeit bestimmt:
\begin{itemize}
	\item die Stärke, Dauer oder Reichweite des Effekts
	\item mögliche Zusatzwirkungen
	\item in manchen Fällen die Kosten (z.\,B. \textit{X -- Rang} AP/MP)
\end{itemize}

Die genaue Skalierung ist bei jeder Fähigkeit einzeln angegeben.

\subsection*{Kosten}

\begin{rulebox}{Kosten durch Wurfzahl}
	Die im Skill- oder Zaubereintrag angegebenen Kosten sind die \textbf{vollen Kosten}.
	Wie viel tatsächlich bezahlt wird, hängt von der \textbf{Wurfzahl} der Ausführung ab.
	Halbe Kosten werden \textbf{aufgerundet}.
\end{rulebox}

\begin{center}
	\begin{tabular}{lp{9cm}}
		\toprule
		\textbf{Wurfzahl} & \textbf{Ergebnis und Kosten} \\
		\midrule
		1--3  & Misslingt, \textbf{halbe Kosten} (aufgerundet) \\
		4--6  & Gelingt, \textbf{volle Kosten} \\
		7--9  & Gelingt, \textbf{halbe Kosten} (aufgerundet) \\
		10+   & Gelingt, \textbf{keine Kosten} \\
		\bottomrule
	\end{tabular}
\end{center}


\subsection*{Legendäre Fähigkeiten}

Fähigkeiten können einen \textbf{legendären Rang (Rang 6)} erreichen.
Dies unterliegt den Regeln des \termref[legendären Rangs]{legendaerer-rang}.

Legendäre Effekte:
\begin{itemize}
	\item sind fest an eine konkrete Fähigkeit gebunden
	\item stehen nur zur Verfügung, wenn diese Fähigkeit Rang 6 erreicht hat
	\item verbrauchen teilweise \termref[Marken]{marken}
\end{itemize}

Die legendären Effekte einer Fähigkeit sind direkt bei der jeweiligen
Fähigkeitsbeschreibung aufgeführt.


\section{Basis-Skills}

Basis-Skills sind nichtmagische Fähigkeiten, die allen Charakteren offenstehen.
Sie verbrauchen in der Regel \textbf{Aktionspunkte (AP)} und werden aktiv
im Kampf oder in Handlungsszenen eingesetzt.

\subsubsection*{Ablenken}
\skillcard
{\term{Ablenken}{ablenken}}
{Bonusaktion}
{4 AP}
{—}
{—}
{
	\begin{itemize}
		\item Ziel erhält einen \textbf{Malus in Höhe des Rangs} auf seine \textbf{nächste \termref[Parade]{parade}}.
	\end{itemize}
}
{
	Der Malus gilt zusätzlich auch für \textbf{\termref[Ausweichen]{ausweichen}}.
}


\subsubsection*{Analyse des Gegners}

\skillcard
{\term{Analyse des Gegners}{analyse-des-gegners}}
{Aktion}
{5 Aktionspunkte (AP)}
{—}
{Rang Runden}
{
	\begin{itemize}
		\item Alle Verbündeten erhalten einen \textbf{+2 Bonus auf Angriffsproben}
		gegen das analysierte Ziel.
	\end{itemize}
}
{
	Die \textbf{aktuellen Lebenspunkte (LP)} des Ziels werden erkannt.
}


\subsubsection*{Erholen}

\skillcard
{\term{Erholen}{erholen}}
{Bonusaktion}
{4 Aktionspunkte (AP)}
{—}
{—}
{
	\begin{itemize}
		\item Stellt \textbf{2 × Rang} Lebenspunkte (LP), Manapunkte (MP)
		\textbf{oder} Aktionspunkte (AP) wieder her.
		\item Ab \textbf{Rang 4} können \textbf{zwei beliebige Reserven}
		gleichzeitig wiederhergestellt werden.
	\end{itemize}
}
{
	Gib \textbf{1 Marke} aus, um einen \textbf{zusätzlichen negativen Effekt}
	vom Charakter zu entfernen.
}


\subsubsection*{Knockdown-Schlag}

\skillcard
{\term{Knockdown-Schlag}{knockdown-schlag}}
{Aktion}
{6 Aktionspunkte (AP)}
{Nahkampf}
{—}
{
	\begin{itemize}
		\item Verursacht \textbf{Waffenschaden + Rang}.
		\item Zusätzlich besteht eine \textbf{1W6-Chance},
		dass das Ziel für \textbf{1 Runde betäubt} wird.
	\end{itemize}
}
{
	Die Chance auf Betäubung wird auf \textbf{1W20} erhöht.
}


\subsubsection*{Reparieren}

\skillcard
{\term{Reparieren}{reparieren}}
{Aktion}
{$5 - \text{Rang}/2$ Aktionspunkte (AP)}
{—}
{—}
{
	\begin{itemize}
		\item Kann nur \textbf{außerhalb des Kampfes} eingesetzt werden.
		\item Repariert \textbf{nichtmagische Gegenstände}
		bis maximal \textbf{80 \% ihres Gesamtwerts}.
		\item \textbf{Magische Gegenstände} können nicht repariert werden.
	\end{itemize}
}
{
	Es besteht eine \textbf{1W4-Chance}, die \textbf{Qualität einer Rüstung}
	(dauerhafte Erhöhung der RP) zu verbessern.
}


\subsubsection*{Schloss knacken}

\skillcard
{\term{Schloss knacken}{schloss-knacken}}
{Aktion}
{2 Aktionspunkte (AP)}
{—}
{—}
{
	\begin{itemize}
		\item Ermöglicht das Öffnen eines \textbf{nichtmagischen Schlosses}
		bis zu einer \textbf{Schlossstufe in Höhe des Rangs}.
	\end{itemize}
}
{
	Gib \textbf{1 Marke} aus, um ein nichtmagisches Schloss
	\textbf{Rang 5 oder niedriger} automatisch zu öffnen.
}


\subsubsection*{Schnelle Ausweichbewegung}

\skillcard
{\term{Schnelle Ausweichbewegung}{schnelle-ausweichbewegung}}
{Bonusaktion}
{2 Aktionspunkte (AP)}
{—}
{—}
{
	\begin{itemize}
		\item Gewährt einen Bonus von \textbf{Rang/2 (aufgerundet)}
		auf das \textbf{nächste Ausweichen}.
		\item Nach einem erfolgreichen Ausweichen darf sich der Charakter
		\textbf{Rang Meter} bewegen.
	\end{itemize}
}
{
	Nach einem erfolgreichen Ausweichen darf der Charakter
	\textbf{sofort eine zusätzliche Bonusaktion} durchführen.
}


\subsubsection*{Schnelle Reflexe}

\skillcard
{\term{Schnelle Reflexe}{schnelle-reflexe}}
{Bonusaktion}
{4 Aktionspunkte (AP)}
{—}
{2 Runden}
{
	\begin{itemize}
		\item Gewährt einen Bonus von \textbf{Rang/2 (aufgerundet)}
		auf \textbf{Ausweichen oder Parade}.
	\end{itemize}
}
{
	Der Bonus kann zusätzlich auf alle
	\textbf{nahen Verbündeten (bis 3 m)} übertragen werden.
}


\subsubsection*{Schneller Ablenkungsschlag}

\skillcard
{\term{Schneller Ablenkungsschlag}{schneller-ablenkungsschlag}}
{Bonusaktion}
{3 Aktionspunkte (AP)}
{Nahkampf (1--2 m)}
{—}
{
	\begin{itemize}
		\item Verursacht \textbf{Schaden in Höhe des Rangs}.
		\item Das Ziel erhält für seine \textbf{nächste Probe}
		einen \textbf{Malus von --1}.
	\end{itemize}
}
{
	Das Ziel \textbf{verliert seine nächste Bonusaktion}.
}


\subsubsection*{Schwachstelle suchen}

\skillcard
{\term{Schwachstelle suchen}{schwachstelle-suchen}}
{Bonusaktion}
{5 Aktionspunkte (AP)}
{—}
{—}
{
	\begin{itemize}
		\item Das Ziel legt eine \textbf{Kampfprobe} gegen deine
		\textbf{Wurfzahl} ab.
		\item Misslingt diese Probe, verursacht der
		\textbf{nächste Angriff} gegen das Ziel
		\textbf{doppelten Schaden}.
	\end{itemize}
}
{
	Gib \textbf{1 Marke} aus, damit der nächste Angriff
	\textbf{Durchschlagsschaden} verursacht.
}


\subsubsection*{Verbände anlegen}

\skillcard
{\term{Verbände anlegen}{verbaende-anlegen}}
{Bonusaktion}
{4 Aktionspunkte (AP)}
{—}
{—}
{
	\begin{itemize}
		\item Heilt \textbf{1W4 + Rang} Lebenspunkte (LP).
	\end{itemize}
}
{
	Gib \textbf{1 Marke} aus, um einen
	\textbf{Statuseffekt} vom Ziel zu entfernen.
}


\subsubsection*{Zähigkeit}

\skillcard
{\term{Zähigkeit}{zaehigkeit}}
{Bonusaktion}
{7 Aktionspunkte (AP)}
{—}
{Ende des Kampfes}
{
	\begin{itemize}
		\item Erhält \textbf{2 × Rang} natürliche
		\textbf{Rüstungspunkte (RP)}.
		\item Kann nur \textbf{einmal pro Kampf} eingesetzt werden.
	\end{itemize}
}
{
	Einmal in diesem Kampf kann der Charakter
	nicht auf \textbf{0 Lebenspunkte} fallen,
	sondern verbleibt stattdessen auf \textbf{1 LP}.
}


\section{Basis-Zauber}

Basis-Zauber sind einfache magische Fähigkeiten, die von allen
zauberkundigen Charakteren erlernt werden können.
Sie verbrauchen in der Regel \textbf{Manapunkte (MP)} und werden aktiv
im Kampf oder in Handlungsszenen eingesetzt.

\subsubsection*{Betören}

\spellcard
{\term{Betören}{betoeren}}
{Bonusaktion}
{8 Manapunkte (MP)}
{5 Meter}
{1W4 + Rang Minuten}
{
	\begin{itemize}
		\item Das Ziel ist dem Zauberer gegenüber \textbf{freundlich}.
		\item Das Ziel kann mit einer \textbf{Wahrnehmungsprobe gegen die Wurfzahl}
		(mit einem Malus von \textbf{–Rang}) widerstehen.
	\end{itemize}
}
{
	Der Zustand \textbf{freundlich} wandelt sich zu \textbf{willenlos}.
	Das Ziel führt nahezu jede Anweisung des Zauberers aus.
}

\subsubsection*{Kleiner Schutzschild}

\spellcard
{\term{Kleiner Schutzschild}{kleiner-schutzschild}}
{Bonusaktion}
{4 Manapunkte (MP)}
{5 Meter}
{Bis bis SP aufgebraucht oder max. 10 Minuten}
{
	\begin{itemize}
		\item Das Ziel erhält \textbf{1 + Rang} \textbf{Magisches Schild (SP)}.
	\end{itemize}
}
{
	Der Zauber kann zusätzlich auf \textbf{einen weiteren Verbündeten}
	angewendet werden.
}


\subsubsection*{Licht erzeugen}

\spellcard
{\term{Licht erzeugen}{licht-erzeugen}}
{Bonusaktion}
{2 Manapunkte (MP)}
{Selbst oder 5 Meter}
{1 Stunde}
{
	\begin{itemize}
		\item Erzeugt eine Lichtkugel, die einen Bereich von
		\textbf{10 Meter + Rang} ausleuchtet.
	\end{itemize}
}
{
	Die Lichtkugel kann in \textbf{Feuer- oder Blitzenergie}
	umgewandelt werden und verursacht
	\textbf{Rang + 3W6 Schaden} in einem Bereich von \textbf{10 Metern}.
}


\subsubsection*{Magischer Schlag}

\spellcard
{\term{Magischer Schlag}{magischer-schlag}}
{Aktion}
{4 Manapunkte (MP)}
{10 Meter}
{—}
{
	\begin{itemize}
		\item Verursacht \textbf{Wurfzahl + Rang Schaden}.
	\end{itemize}
}
{
	Es können \textbf{zwei Geschosse} erzeugt werden,
	beide mit vollem Schaden.
}


\subsubsection*{Magiesinn}

\spellcard
{\term{Magiesinn}{magiesinn}}
{Bonusaktion}
{2 Manapunkte (MP)}
{Selbst}
{10 Minuten}
{
	\begin{itemize}
		\item Magie kann in einem Radius von \textbf{10 + Rang Metern} wahrgenommen werden.
		\item Es wird ausschließlich das \textbf{Vorhandensein} von Magie erkannt,
		keine Details.
	\end{itemize}
}
{
	Der Zauber erkennt zusätzlich \textbf{Art, Stärke und Quelle}
	magischer Energien.
}


\subsubsection*{Magisches Schloss öffnen}

\spellcard
{\term{Magisches Schloss öffnen}{magisches-schloss-oeffnen}}
{Bonusaktion}
{4 Manapunkte (MP)}
{5 Meter}
{—}
{
	\begin{itemize}
		\item Öffnet ein \textbf{nichtmagisches Schloss} oder eine vergleichbare Mechanik.
		\item Bei \textbf{magischen Schlössern, Siegeln oder Barrieren} ist zusätzlich
		eine erfolgreiche \textbf{Wissensprobe} erforderlich.
	\end{itemize}
}
{
	Der Zauberer erkennt automatisch das \textbf{Vorhandensein}
	und den \textbf{Rang} eines Schlosses, einer Verriegelung
	oder eines magischen Siegels.
}


\subsubsection*{Mana aufladen}

\spellcard
{\term{Mana aufladen}{mana-aufladen}}
{Bonusaktion}
{2 Manapunkte (MP)}
{Selbst}
{—}
{
	\begin{itemize}
		\item Stellt \textbf{Rang + 1W8 Manapunkte (MP)} wieder her.
		\item Kann nur \textbf{einmal pro Stunde} eingesetzt werden.
		\item Die \textbf{Zauberdauer beträgt 5 Minuten}.
	\end{itemize}
}
{
	Die Zauberdauer entfällt, Mana wird sofort wiederhergestellt.
	Zusätzlich kann der Zauber auf \textbf{einen Verbündeten}
	angewandt werden.
}


\subsubsection*{Trickillusion}

\spellcard
{\term{Trickillusion}{trickillusion}}
{Bonusaktion}
{2 + Rang Manapunkte (MP)}
{10 Meter}
{10 Minuten}
{
	\begin{itemize}
		\item Erzeugt eine visuelle oder akustische Illusion von bis zu
		\textbf{1 + Rang Quadratmetern}.
		\item Die Illusion ist unbeweglich.
		\item Sie kann durch eine \textbf{Wahrnehmungsprobe} durchschaut werden.
	\end{itemize}
}
{
	Illusionen erhalten einen \textbf{physischen Körper} und können
	kleine Aktionen durchführen.
	Angriffe beenden die Illusion sofort.
}


\subsubsection*{Unsichtbare Hand}

\spellcard
{\term{Unsichtbare Hand}{unsichtbare-hand}}
{Aktion}
{2 Manapunkte (MP)}
{10 Meter + Rang}
{1 Minute}
{
	\begin{itemize}
		\item Die Hand kann Objekte bis zu \textbf{5 + Rang kg} bewegen oder manipulieren.
		\item Sie kann \textbf{keine Angriffe} ausführen.
	\end{itemize}
}
{
	Die Hand kann nun direkte Angriffe durchführen.
	Der verursachte Schaden wird mit dem Spielleiter abgestimmt.
}


\subsubsection*{Unsichtbarkeit}

\spellcard
{\term{Unsichtbarkeit}{unsichtbarkeit}}
{Bonusaktion}
{6 Manapunkte (MP)}
{Selbst (ab Rang 4 zusätzlich 1 Verbündeter innerhalb von 2 m)}
{Rang Runden oder bis gebrochen}
{
	\begin{itemize}
		\item Ziel wird unsichtbar.
		\item Nichtmagische Gegner können das Ziel nur über eine
		\textbf{Wahrnehmungsprobe gegen Wurfzahl + Rang} lokalisieren.
		\item \textbf{Magiesinn} erkennt lediglich das Vorhandensein von Magie,
		keine genaue Position.
		\item Angriffe oder Schadenzauber brechen den Effekt sofort.
		\item Während Unsichtbarkeit dürfen nur Zauber mit
		\textbf{Manakosten $\le 5$} gewirkt werden.
	\end{itemize}
}
{
	Gib \textbf{1 Marke} aus, um einen Angriff oder Zauber beliebiger Kosten
	auszuführen, ohne den Effekt sofort zu brechen.
	Der Zauber endet erst am Ende der nächsten Runde.
}


\subsubsection*{Vergessen}

\spellcard
{\term{Vergessen}{vergessen}}
{Aktion}
{8 Manapunkte (MP)}
{5 Meter}
{—}
{
	\begin{itemize}
		\item Das Ziel vergisst bis zu \textbf{Rang $\times$ 10 Minuten}
		vergangener Handlungen.
		\item Widerstand über eine \textbf{Wahrnehmungsprobe gegen die Wurfzahl}.
	\end{itemize}
}
{
	Erinnerungen, Zauber, Skills oder Talente können
	\textbf{dauerhaft} vergessen werden.
	Jeder dauerhafte Einsatz erhöht den Widerstand gegen weiteres Vergessen
	um \textbf{+3}.
}


\subsubsection*{Wunde schließen}

\spellcard
{\term{Wunde schließen}{wunde-schliessen}}
{Aktion}
{4 Manapunkte (MP)}
{Berührung}
{—}
{
	\begin{itemize}
		\item Heilt \textbf{1W6 + Rang Lebenspunkte (LP)}.
	\end{itemize}
}
{
	Der Zauberer heilt sich selbst zusätzlich um
	\textbf{Rang Lebenspunkte (LP)}.
}


\subsubsection*{Zaubermanipulation}

\spellcard
{\term{Zaubermanipulation}{zaubermanipulation}}
{Bonusaktion}
{$4 + \text{Rang}$ Manapunkte (MP)}
{15 m (fremde Zauber), Selbst (eigene Zauber)}
{—}
{
	Zaubermanipulation kann auf \textbf{eigene} oder \textbf{fremde} Zauber angewendet werden:
	\begin{itemize}
		\item \textbf{Eigene Zauber:} Manipulation als Bonusaktion beim Wirken
		oder direkt im Anschluss.
		\item \textbf{Fremde Zauber:} Manipulation erfolgt \textbf{sofort};
		die eigene Bonusaktion der Runde wird vorgezogen.
	\end{itemize}
	
	\vspace{4pt}
	\textbf{Effekte je Rang:}
	\begin{itemize}
		\item \textbf{Rang 1:} Änderung des Schadenstyps für \textbf{Rang} und
		\textbf{2 Runden} oder kleiner Parameter.
		\item \textbf{Rang 2:} Änderung von Stärke oder Form.
		\item \textbf{Rang 3:} Hinzufügen eines Zusatzeffekts.
		\item \textbf{Rang 4:} Tiefgreifende Umstrukturierung (z.\,B. Umkehr).
		\item \textbf{Rang 5:} Komplexe Umstrukturierung ohne Nebenwirkungen.
	\end{itemize}
}
{
	Es können \textbf{zwei Effekte gleichzeitig} ausgeführt werden.
}

\section{Klassen (nicht magisch)}

\term{Nicht Magische Klassen}{nicht-magische-klassen} haben jeweils 1 Startausrüstung.

\subsection{Krieger}

\subsubsection*{Rolle}
Tank / Nahkämpfer

\subsubsection*{Kurzprofil}
Der \term{Krieger}{klasse-krieger} ist der Frontkämpfer mit schwerer Rüstung. Der Krieger bindet Gegner, schützt
Verbündete und übt konstanten Druck im Nahkampf aus.

\subsubsection*{Stärken}
\begin{itemize}
	\item Hohe Lebenspunkte
	\item Starke defensive Fähigkeiten
	\item Vielseitige Waffenwahl
	\item Kontrolle über Gegner durch Provozieren und Flächeneffekte
\end{itemize}

\subsubsection*{Kampfstil}
Direkt, offensiv, standhaft – ideal für die erste Reihe.

% =========================
% Startausrüstung
% =========================
\subsubsection*{Startausrüstung}

\begin{center}
	\begin{tabular}{lp{9cm}}
		\toprule
		\textbf{Gegenstand} & \textbf{Wert und Eigenschaften} \\
		\midrule
		Schwere Rüstung (normal) & 4 Rüstungspunkte (RP), --1 auf \textbf{Gewandtheit} \\
		Schild (normal) & 2 Rüstungspunkte (RP) \\
		Waffe nach Wahl &
		Kurzschwert (\textbf{Stärke} +2 Schaden) \\
		& oder Streitaxt (\textbf{Stärke} +3 Schaden, --1 auf Angriffswürfe) \\
		\bottomrule
	\end{tabular}
\end{center}

% =========================
% Klassenfähigkeiten
% =========================
\subsubsection*{Klassenfähigkeiten}

\paragraph*{Wuchtschlag}

\skillcard
{\term{Wuchtschlag}{wuchtschlag}}
{Bonusaktion}
{$2 + \text{Rang}$ AP}
{Nahkampf}
{—}
{
	\begin{itemize}
		\item Verursacht \textbf{3 + Rang Durchschlagsschaden}.
		\item Schaden ignoriert Rüstung.
	\end{itemize}
}
{
	Wähle zusätzlich \textbf{Verwirrung} oder \textbf{Erschöpfung} als Statuseffekt für das Ziel.
}

\paragraph*{Schildwall}
\skillcard
{\term{Schildwall}{schildwall}}
{Aktion}
{$4 + \text{Rang}$ AP}
{Selbst, Radius 3 m}
{1 Runde}
{
	\begin{itemize}
		\item Der Krieger und bis zu \textbf{Rang Verbündete} erleiden
		\textbf{Wurfzahl} weniger \textbf{normalen Schaden}.
		\item Gilt nicht für Durchschlagsschaden.
	\end{itemize}
}
{
	Der Effekt reduziert nun auch \textbf{Durchschlagsschaden}.
}

\paragraph*{Provozieren}
\skillcard
{\term{Provozieren}{provozieren}}
{Bonusaktion}
{4 AP}
{Bis zu 10 m}
{1 Runde}
{
	\begin{itemize}
		\item Zwingt Gegner, den Krieger anzugreifen.
		\item Effekte skalieren mit dem Rang:
	\end{itemize}
	
	\begin{center}
		\begin{tabular}{cccc}
			\toprule
			\textbf{Rang} & \textbf{Ziele} & \textbf{Reichweite} & \textbf{Zusatz} \\
			\midrule
			1 & 1 & 4 m & Ziel greift nur den Krieger an \\
			2 & 2 & 6 m & +1 RP gegen diese Ziele \\
			3 & 3 & 8 m & Gegner: --1 Angriff, Krieger: +2 RP \\
			4 & 4 & 10 m & Dornen (2 Schaden bei Nahkampf) \\
			5 & 5 & 10 m & Verbündete: +1 auf alle Proben \\
			\bottomrule
		\end{tabular}
	\end{center}
}
{
	Für jeden Gegner, der den Krieger angreift, erhält er
	\textbf{eine zusätzliche Bonusaktion} für den Rest des Kampfes.
}

\paragraph*{Standhaftigkeit}
\skillcard
{\term{Standhaftigkeit}{standhaftigkeit}}
{Bonusaktion}
{$3 + \text{Rang}$ AP}
{Selbst}
{Bis Kampfende}
{
	\begin{itemize}
		\item Solange LP $\ge 50\,\%$: \textbf{+Rang/2} Bonus auf Verteidigungsproben.
		\item Solange LP $< 50\,\%$: \textbf{+Rang} Bonus auf Verteidigungsproben.
	\end{itemize}
}
{
	Alle Schadenstypen zählen für den Krieger als \textbf{normaler Schaden}.
}

\paragraph*{Klingenwirbel}
\skillcard
{\term{Klingenwirbel}{klingenwirbel}}
{Aktion}
{8 AP}
{Nahkampf (Umkreis)}
{—}
{
	\begin{center}
		\begin{tabular}{ccc}
			\toprule
			\textbf{Rang} & \textbf{Ziele} & \textbf{Schaden} \\
			\midrule
			1 & 2 & 1 + 1W4 \\
			2 & 2 & 1 + 1W6 \\
			3 & 3 & 2 + 1W8 \\
			4 & 3 & 3 + 1W8 \\
			5 & Alle & 3 + 1W10 \\
			6 & Alle & 4 + 1W12 \\
			\bottomrule
		\end{tabular}
	\end{center}
}
{
	Der gesamte Schaden wird zu \textbf{Durchschlagsschaden}.
}

\paragraph*{(fehlender 6. Zauber)}
% TODO:  6. Klassenzauber liefern

\subsection{Schurke}

\subsubsection*{Rolle}
Assassine / Heimlichkeit / Kontrolle

\subsubsection*{Kurzprofil}
Der \term{Schurke}{klasse-schurke} ist ein Spezialist für Hinterhalte, Präzision und das Ausnutzen von Schwächen.
Er arbeitet aus dem Schatten, manipuliert Situationen und entscheidet Kämpfe mit gezielten Angriffen.

\subsubsection*{Stärken}
\begin{itemize}
	\item Hohe Mobilität und starke Positionierung
	\item Hoher Burst-Schaden durch Hinterhalte
	\item Werkzeuge für Kontrolle, Täuschung und Infiltration
	\item Gut im Umgang mit Fallen und Schlössern
\end{itemize}

\subsubsection*{Kampfstil}
Schnell, opportunistisch, tödlich – ideal für Überraschungsangriffe und das Ausschalten einzelner Ziele.

% =========================
% Startausrüstung
% =========================
\subsubsection*{Startausrüstung}

\begin{center}
	\begin{tabular}{lp{9cm}}
		\toprule
		\textbf{Gegenstand} & \textbf{Wert und Eigenschaften} \\
		\midrule
		Umhang & +1 Bonus auf \textbf{HEIMLICHKEIT}-Wurf \\
		Mittlere Rüstung (normal) & 3 Rüstungspunkte (RP) \\
		Dolch & \textbf{STÄRKE} +1 Schaden \\
		Kurzschwert & \textbf{STÄRKE} +2 Schaden \\
		\bottomrule
	\end{tabular}
\end{center}

% =========================
% Klassenfähigkeiten
% =========================
\subsubsection*{Klassenfähigkeiten}

\paragraph*{Hinterhalt}
\skillcard
{\term{Hinterhalt}{schurke-hinterhalt}}
{Bonusaktion}
{3 Aktionspunkte (AP)}
{Nahkampf}
{Bis zum nächsten Angriff}
{
	\begin{itemize}
		\item Verdoppelt den Schaden des \textbf{nächsten Angriffs}.
		\item Nur einsetzbar bei \textbf{unvorbereiteten Zielen}.
		\item Kann nur verwendet werden, wenn der Gegner den Schurken
		\textbf{nicht im Fokus} hat.
	\end{itemize}
}
{
	Gib \textbf{1 Marke} aus, um dem Ziel \textbf{Vergiftung, Verbrennung}
	oder \textbf{Unterkühlung} zuzufügen.
}

\paragraph*{Manipulation}
\skillcard
{\term{Manipulation}{schurke-manipulation}}
{Aktion}
{3 Aktionspunkte (AP)}
{—}
{—}
{
	\begin{itemize}
		\item Wenn die \textbf{Wahrnehmungsprobe} des Ziels gegen die
		\textbf{Wurfzahl} des Schurken misslingt, darf eine Manipulation
		durchgeführt werden (z.\,B. Überreden, Ablenken, Täuschen).
	\end{itemize}
}
{
	Gib \textbf{1 Marke} aus, um die \textbf{Schwachstelle} des Gegenübers
	zu erfahren. Das Ausnutzen dieser Schwachstelle gewährt
	\textbf{+3 Bonus} auf die Manipulation.
}

\paragraph*{Feinmotorik}
\skillcard
{\term{Feinmotorik}{schurke-feinmotorik}}
{Aktion}
{$2 + \text{Rang des Ziels} - \text{Rang}$ Aktionspunkte (AP)}
{—}
{—}
{
	\begin{itemize}
		\item \textbf{Fallenentschärfen:} Erkennt und deaktiviert Fallen
		automatisch bei erfolgreicher Probe.
		\item \textbf{Taschendiebstahl:} Entwendet einen kleinen Gegenstand.
		Misslingt die Wahrnehmungsprobe des Ziels gegen die Wurfzahl des
		Schurken, bleibt der Diebstahl unbemerkt.
	\end{itemize}
}
{
	Erlaubt das Nutzen der Effekte auf \textbf{Distanz bis 2 m}
	(z.\,B. Draht, Haken, Spiegel).
}

\paragraph*{Giftklinge}
\skillcard
{\term{Giftklinge}{schurke-giftklinge}}
{Bonusaktion}
{3 Aktionspunkte (AP)}
{Selbst (Waffe)}
{1 Treffer oder 2 Runden}
{
	\begin{itemize}
		\item Die Waffe wird vergiftet.
		\item Zusätzlich zum Waffenschaden verursacht jeder Treffer
		\textbf{Rang 1 Giftschaden für 2 Runden}.
		\item Ab \textbf{Rang 3}: \textbf{Rang 2 Giftschaden}.
		\item Ab \textbf{Rang 5}: \textbf{Rang 3 Giftschaden}.
	\end{itemize}
}
{
	Gib \textbf{1 Marke} aus, damit der nächste Treffer zusätzlich
	\textbf{Erosion Rang 3} verursacht.
}

\paragraph*{Mordlust}
\skillcard
{\term{Mordlust}{schurke-mordlust}}
{Aktion}
{6 Aktionspunkte (AP)}
{2 + Rang Meter}
{Bis Ziel besiegt oder geflohen}
{
	\begin{itemize}
		\item Zusätzlich zum Waffenschaden:
		Jeder Treffer verursacht Giftschaden:
		\begin{itemize}
			\item Rang 1--2: \textbf{Rang 1} für 2 Runden
			\item Rang 3--4: \textbf{Rang 2} für 2 Runden
			\item Rang 5: \textbf{Rang 3} für 2 Runden
		\end{itemize}
		\item Die \textbf{Talente des Ziels} werden für \textbf{1 Runde halbiert}.
	\end{itemize}
	
	\textbf{Negativer Trade:}
	\begin{itemize}
		\item Verbündete müssen eine \textbf{Wahrnehmungsprobe} bestehen,
		ansonsten verlieren sie ihre \textbf{Bonusaktion}.
	\end{itemize}
}
{
	Das Ausschalten des Ziels durch den Schurken
	\textbf{resettet Aktion und Bonusaktion} in dieser Runde.
}

\paragraph*{Schattenschritt}
\skillcard
{\term{Schattenschritt}{schurke-schattenschritt}}
{Aktion}
{6 Aktionspunkte (AP)}
{Sichtbarer Schattenbereich}
{—}
{
	\begin{itemize}
		\item Bewegt sich bis zu \textbf{10 Meter} in einen sichtbaren
		schattigen Bereich.
		\item Der nächste Angriff erhält einen \textbf{Rang Bonus}
		auf den Angriffswurf.
		\item Kann genutzt werden, um Hindernisse wie Mauern oder
		feindliche Gruppen zu umgehen.
	\end{itemize}
}
{
	Nach jedem Schattenschritt ist der Schurke bis zur
	nächsten Runde \textbf{verborgen}.
	Greift er an, endet die Verborgenheit nach dem Angriff.
}

\subsection{Waldläufer}

\subsubsection*{Rolle}
Fernkämpfer / Überlebenskünstler

\subsubsection*{Kurzprofil}
Der \term{Waldläufer}{klasse-waldlaeufer} ist ein Fernkampf-Spezialist mit starker
Naturverbundenheit. Er hält Gegner auf Distanz, erkennt Gefahren frühzeitig und nutzt
Tiere sowie Gelände zu seinem Vorteil.

\subsubsection*{Stärken}
\begin{itemize}
	\item Effektiv im Fernkampf
	\item Naturkenntnis (Spuren, Gefahren, Ressourcen)
	\item Vielseitige Tierunterstützung (Schaden, Kontrolle, Buffs)
\end{itemize}

\subsubsection*{Kampfstil}
Präzise Angriffe aus sicherer Entfernung, taktisch und umsichtig.

\subsubsection*{Startausrüstung}

\begin{center}
	\begin{tabular}{lp{9cm}}
		\toprule
		\textbf{Gegenstand} & \textbf{Wert und Eigenschaften} \\
		\midrule
		Leichte Rüstung (normal) & 2 Rüstungspunkte (RP) \\
		Reflexbogen & 1W6 Schaden \\
		Kurzschwert & Stärke +2 Schaden \\
		Heilkräuter (1×) & Heilt 1W4 Lebenspunkte (LP) \\
		Materialien (5×) & Nutzbar als Nahrung oder kreative Einfälle \\
		\bottomrule
	\end{tabular}
\end{center}

Der Waldläufer startet jedes Abenteuer mit 5 Materialien.


\begin{rulebox}{Hinweis: Pfeile}
	Pfeile zerbersten bei:
	\begin{itemize}
		\item \textbf{+15 Schaden} allgemein
		\item \textbf{+13 Schaden} auf harten Oberflächen (Rüstung, Stein, Eisen)
	\end{itemize}
\end{rulebox}

\subsubsection*{Klassenfähigkeiten}

\paragraph*{Gezielter Schuss}
\skillcard
{\term{Gezielter Schuss}{waldlaeufer-gezielter-schuss}}
{Aktion}
{6 Aktionspunkte (AP)}
{Fernkampf}
{—}
{
	\begin{itemize}
		\item Verursacht zusätzlich zum Waffenschaden
		\textbf{Wurfzahl Durchschlagsschaden}.
	\end{itemize}
}
{
	Unter Aufgabe der Bonusaktion dieser Runde
	kann ein \textbf{Zusatzeffekt} hinzugefügt werden.
}

\paragraph*{Waldkundiger Blick}
\skillcard
{\term{Waldkundiger Blick}{waldlaeufer-waldkundiger-blick}}
{Bonusaktion}
{2 Aktionspunkte (AP)}
{Selbst}
{Modusabhängig}
{
	Bei Aktivierung wird ein Modus gewählt:
	
	\textbf{Modus A – Spuren \& Gefahren (1 Minute):}
	\begin{itemize}
		\item +2 Bonus auf Wahrnehmung und Wildnisleben
		(+3 ab Rang 3).
		\item Deckt im Radius von \textbf{5 m × Rang} versteckte Kreaturen
		oder Fallen auf.
		\item Bei Wurfzahl 10+: Überraschungen gegen dich oder deine Gruppe
		werden negiert.
	\end{itemize}
	
	\textbf{Modus B – Kräuter \& Ressourcen (10 Minuten):}
	\begin{itemize}
		\item 4–5: 1× Heilkräuter
		\item 6–9: 2× Heil-, Stärkungs- oder Fokuskraut
		\item 10+: Rang Bündel (freie Mischung)
		\item Jeder Fund kann als Materialien gewertet werden
		(Verbrauch der Pflanze).
	\end{itemize}
	
	\textbf{Modus C – Gelände-Vorteil (1 Runde):}
	\begin{itemize}
		\item Markiere Ziel oder Zone (2 m Radius).
		\item +1 Bonus auf nächste Angriffsprobe
		\item +1 Bonus auf Ausweichen in der markierten Zone.
	\end{itemize}
}
{
	Beim Betreten von Natur- oder Außengelände gibt die SL
	einen klaren Hinweis (Route, Gefahrenzone oder Engstelle).
}

\paragraph*{Natürlicher Schutz}
\skillcard
{\term{Natürlicher Schutz}{waldlaeufer-natuerlicher-schutz}}
{Bonusaktion}
{4 Aktionspunkte (AP)}
{Selbst oder Verbündeter}
{2 Runden}
{
	\begin{itemize}
		\item +2 Bonus auf Verteidigungsproben (Ausweichen oder Parade).
		\item Umgebung muss nutzbare Deckung bieten
		(Entscheidung der SL).
	\end{itemize}
}
{
	Materialien können genutzt werden,
	um aktiv eine Deckung zu errichten.
}

\paragraph*{Mehrfachschuss}
\skillcard
{\term{Mehrfachschuss}{waldlaeufer-mehrfachschuss}}
{Aktion}
{4 Aktionspunkte (AP)}
{Fernkampf}
{—}
{
	\begin{itemize}
		\item Feuert \textbf{Rang + 1 Pfeile} auf gewählte Ziele.
		\item Jeder Treffer verursacht \textbf{halben Waffenschaden}.
	\end{itemize}
}
{
	Der erste Pfeil verursacht \textbf{vollen Waffenschaden},
	alle weiteren weiterhin halben Schaden.
}

\paragraph**{Tiergefährte rufen}
\skillcard
{\term{Tiergefährte rufen}{tiergefaehrte-rufen}}
{Aktion}
{$5 - \text{Rang}$ Aktionspunkte (AP) + 1 Materialien}
{Selbst}
{—}
{
	\begin{itemize}
		\item Kann nur \textbf{einmal pro Kampf} eingesetzt werden.
		\item Ruft einen Tiergefährten, der den Waldläufer unterstützt.
		\item Der Tiergefährte bleibt für \textbf{2 Runden} aktiv.
	\end{itemize}
	
	\medskip
	\textbf{Standard-Tiergefährten (eines wählbar):}
	
	\begin{itemize}
		\item \textbf{Fuchs:}
		Greift für \textbf{2 Runden} an.  
		Jeder Angriff verursacht \textbf{1W6 Schaden}.
		
		\item \textbf{Schwalbe:}
		Fliegt für \textbf{2 Runden}.  
		Markiert ein Ziel: Alle Verbündeten erhalten für \textbf{1 Runde}
		\textbf{+1 Bonus auf Angriffsproben} gegen dieses Ziel.
		
		\item \textbf{Eichhörnchen:}
		Stiehlt oder beschädigt Ausrüstung.  
		Ziel erhält \textbf{–1 auf Angriffsproben} für \textbf{2 Runden}.
		
		\item \textbf{Dachs:}
		Ignoriert \textbf{Rüstung (RP)}.  
		Verursacht \textbf{1W4 Durchschlagsschaden} pro Runde für \textbf{2 Runden}.  
		Ziel erhält \textbf{–2 auf Gewandtheitsproben}.
	\end{itemize}
}
{
	\textbf{Variante A – Geistergefährte}
	\begin{itemize}
		\item Der Geistergefährte bleibt für \textbf{4 Runden}.
		\item In jeder Runde kann er sich für \textbf{eine Aktion}
		in eines der \textbf{Standardtiere} manifestieren.
	\end{itemize}
	
	\medskip
	\textbf{Variante B – Legendäres Tier (eines wählbar):}
	
	\medskip
	\textbf{Schattenwolf}
	\begin{itemize}
		\item Greift für \textbf{3 Runden} an.
		\item Jeder Angriff verursacht:
		\begin{itemize}
			\item \textbf{1W8 Giftschaden}
			\item \textbf{Gift Rang 3} für \textbf{3 Runden}
		\end{itemize}
		\item Ziel gilt als \textbf{markiert}.
		\item Alle Verbündeten erhalten bis zum Ende der nächsten Runde
		\textbf{+1 Bonus auf ihre nächste Angriffsprobe}
		gegen das markierte Ziel.
	\end{itemize}
	
	\medskip
	\textbf{Himmelsfalke}
	\begin{itemize}
		\item Fliegt für \textbf{1 Runde}.
		\item In Runde 2: \textbf{Sturzflug}.
		\item Markiert \textbf{2 Ziele}:
		\begin{itemize}
			\item Diese verlieren für \textbf{2 Runden ihre Aktion}.
			\item Ein Ziel erleidet \textbf{Blendung Rang 3}.
		\end{itemize}
	\end{itemize}
	
	\medskip
	\textbf{Schwarzbär}
	\begin{itemize}
		\item Greift für \textbf{3 Runden} an.
		\item Jeder Angriff verursacht \textbf{1W10 Durchschlagsschaden}.
		\item Absorbiert bis zu \textbf{40 Schaden},
		der eigentlich auf Verbündete gehen würde.
	\end{itemize}
}

\paragraph*{(fehlender 6. Zauber)}
% TODO:  6. Klassenzauber liefern

\subsection{Gladiator}

\subsubsection*{Rolle}
Schadensausteiler mit Stil

\subsubsection*{Kurzprofil}
Der \term{Gladiator}{klasse-gladiator} ist ein charismatischer Nahkämpfer, der
physischen Schaden mit Präsenz und Auftritt verbindet.
Er beeinflusst Gegner wie Verbündete durch spektakuläre Angriffe und Motivation.

\subsubsection*{Stärken}
\begin{itemize}
	\item Hoher Nahkampfschaden
	\item Flächeneffekte und Mobilität
	\item Starke soziale Präsenz (Einschüchterung, Motivation)
\end{itemize}

\subsubsection*{Kampfstil}
Stilvoll, aggressiv, publikumswirksam.

% =========================
% Startausrüstung
% =========================
\subsubsection*{Startausrüstung}

\begin{center}
	\begin{tabular}{lp{9cm}}
		\toprule
		\textbf{Gegenstand} & \textbf{Wert und Eigenschaften} \\
		\midrule
		Mittlere Rüstung (normal) & 3 Rüstungspunkte (RP) \\
		Kurzschwert & \textbf{Stärke} +2 Schaden \\
		Große Waffe nach Wahl &
		Langschwert (2× \textbf{Stärke} Schaden) \\
		& oder Kriegsaxt (2× \textbf{Stärke} +3 Schaden, zerstört 2 RP, --2 auf Verteidigungsproben) \\
		\bottomrule
	\end{tabular}
\end{center}

% =========================
% Klassenfähigkeiten
% =========================
\subsubsection*{Klassenfähigkeiten}

\paragraph*{Publikumswirksamkeit}

\skillcard
{\term{Publikumswirksamkeit}{publikumswirksamkeit}}
{Bonusaktion}
{2 AP}
{Radius \textbf{Rang} m}
{—}
{
	\begin{itemize}
		\item Verbündete im Radius erhalten \textbf{Rang/2 (aufgerundet)}
		Bonus auf ihre \textbf{nächste Angriffsprobe}.
		\item Ab \textbf{Rang 3}: Gegner im Radius müssen eine
		\textbf{Stärke-Probe} ablegen.
		\item Bei Misserfolg: \textbf{--1 Malus} auf ihre nächste Angriffsprobe.
	\end{itemize}
}
{
	Gegner erhalten zusätzlich \textbf{Verwirrung (1 Runde)}.
	Verbündete erhalten \textbf{1 zusätzliche Bonusaktion}.
}

\paragraph*{Einschüchternder Schlag}

\skillcard
{\term{Einschüchternder Schlag}{einschuechternder-schlag}}
{Aktion}
{4 AP}
{Nahkampf}
{—}
{
	\begin{itemize}
		\item Verursacht \textbf{Waffenschaden + Rang}.
		\item Ziel legt eine \textbf{Kampf-Probe} ab.
		\item Bei Misserfolg: \textbf{--2 Malus auf alle Proben} für \textbf{1 Runde}.
	\end{itemize}
}
{
	Gib \textbf{1 Marke} aus:
	Beim nächsten Treffer weicht das Ziel \textbf{3 m zurück}
	und verliert seine \textbf{Bonusaktion für 2 Runden}.
}

\paragraph*{Tanzender Kampf}

\skillcard
{\term{Tanzender Kampf}{tanzender-kampf}}
{Aktion}
{4 AP}
{Nahkampf}
{—}
{
	\begin{itemize}
		\item Führt einen Angriff mit \textbf{normalem Waffenschaden} aus.
		\item Danach Bewegung um bis zu \textbf{2 + Rang Meter}.
	\end{itemize}
}
{
	Fernkampfangriffe gegen den Gladiator erhalten
	\textbf{--3 Malus} für \textbf{2 Runden}.
}

\paragraph*{Klingensturm}

\skillcard
{\term{Klingensturm}{klingensturm-gladiator}}
{Aktion}
{5 AP}
{Nahbereich (3 m)}
{—}
{
	\begin{itemize}
		\item Trifft \textbf{alle Gegner} im Nahbereich.
		\item Jeder erleidet \textbf{3 + Rang Schaden}.
		\item Für jeden getroffenen Gegner erhält der Gladiator
		\textbf{+1 Bonus auf die nächste Probe}.
	\end{itemize}
}
{
	Jeder Einsatz von Klingensturm in diesem Kampf reduziert
	die AP-Kosten zukünftiger Einsätze um \textbf{1}.
}

\paragraph*{Inspirierendes Brüllen}

\skillcard
{\term{Inspirierendes Brüllen}{inspirierendes-bruellen}}
{Aktion}
{4 AP}
{—}
{—}
{
	\begin{itemize}
		\item Kann nur \textbf{einmal pro Kampf} eingesetzt werden.
	\end{itemize}
	
	\medskip
	\textbf{Effekt nach Rang:}
	
	\begin{center}
		\begin{tabular}{cl}
			\toprule
			\textbf{Rang} & \textbf{Effekt} \\
			\midrule
			1 & Verbündete im Umkreis von 3 m: +1 \textbf{Kampf} für 1 Runde \\
			2 & Verbündete im Umkreis von 5 m: +1 \textbf{Kampf} für 2 Runden \\
			3 & Alle Verbündeten: +1 \textbf{Kampf} für diesen Kampf \\
			4 & Alle Verbündeten: +1 \textbf{Kampf} und \textbf{Stärke} \\
			5 & Alle Verbündeten und der Gladiator:
			+2 \textbf{Kampf} und \textbf{Stärke} \\
			\bottomrule
		\end{tabular}
	\end{center}
}
{
	Gib \textbf{2 Marken} aus:
	Wähle einen Zusatzeffekt und befreie alle Verbündeten
	von diesem negativen Zustand.
}

\paragraph*{(fehlender 6. Zauber)}
% TODO:  6. Klassenzauber liefern

\subsection{Mönch}

\subsubsection*{Rolle}
Beweglicher Nahkämpfer / Spiritueller Unterstützer

\subsubsection*{Kurzprofil}
Der \term{Mönch}{klasse-moenich} ist ein agiler Nahkämpfer, der Körper und Geist
durch spirituelle Disziplin perfektioniert.
Er verzichtet auf schwere Rüstung und Waffen und nutzt stattdessen
Beweglichkeit, präzise Schläge und spirituelle Kräfte.

\subsubsection*{Stärken}
\begin{itemize}
	\item Hohe Mobilität und Ausweichfähigkeit
	\item Effektiv ohne schwere Rüstung
	\item Mischung aus Schaden, Kontrolle und Heilung
	\item Teilweise unabhängig von externen Heilern
\end{itemize}

\subsubsection*{Kampfstil}
Präzise Schläge, Ausweichmanöver und spirituelle Techniken.

% =========================
% Startausrüstung
% =========================
\subsubsection*{Startausrüstung}

\begin{center}
	\begin{tabular}{lp{9cm}}
		\toprule
		\textbf{Gegenstand} & \textbf{Wert und Eigenschaften} \\
		\midrule
		Mönchskutte & +1 Bonus auf \textbf{GEWANDTHEIT-Wurfzahl} \\
		Kampfhandschuhe & +1 Bonus auf \textbf{STÄRKE-Wurfzahl} \\
		Meditationsperlen & Nach 5 Minuten Meditation: +2 MP \\
		\bottomrule
	\end{tabular}
\end{center}

% =========================
% Klassenfähigkeiten
% =========================
\subsubsection*{Klassenfähigkeiten}

\paragraph*{Schlagserie}

\skillcard
{\term{Schlagserie}{schlagserie}}
{Aktion}
{5 AP}
{Nahkampf}
{—}
{
	\begin{itemize}
		\item Führt \textbf{zwei Angriffe} aus.
		\item Jeder Angriff verursacht \textbf{halben Schaden}:
		\[
		\frac{\text{STÄRKE-Wert} + \text{Wurfzahl}}{2} \ (\text{abgerundet})
		\]
		\item Angriffe können auf \textbf{ein oder zwei Ziele} verteilt werden.
	\end{itemize}
}
{
	Ein zusätzlicher \textbf{dritter finaler Schlag}:
	Schaden = \textbf{STÄRKE-Wert + Wurfzahl}.
}

\paragraph*{Konzentrierter Schlag}

\skillcard
{\term{Konzentrierter Schlag}{konzentrierter-schlag}}
{Aktion}
{8 AP}
{Nahkampf}
{2 Runden}
{
	\begin{itemize}
		\item \textbf{Runde 1:} Vorbereitung —
		\textbf{--1} auf Ausweichen oder Parade.
		\item \textbf{Runde 2:} Angriff verursacht
		\textbf{Rang × STÄRKE-Wert Schaden}.
	\end{itemize}
}
{
	\begin{itemize}
		\item Schaden wird \textbf{1:1 von Rüstung (RP)} abgezogen.
		\item Ziel erhält \textbf{Erschöpfung} oder verliert seine \textbf{Bonusaktion}.
	\end{itemize}
}

\paragraph*{Wirbeltritt}

\skillcard
{\term{Wirbeltritt}{wirbeltritt}}
{Aktion}
{5 AP}
{Nahbereich (2 m)}
{—}
{
	\begin{itemize}
		\item Trifft alle Gegner im Nahbereich.
		\item Jeder erleidet \textbf{STÄRKE-Wert + Rang Schaden}.
		\item Jeder Gegner legt eine \textbf{Stärke-Probe} ab.
		\item Bei Misserfolg: \textbf{1 m Rückstoß}.
	\end{itemize}
}
{
	Der Schaden gilt als \textbf{Direktschaden}.
}

\paragraph*{Meditation zur Heilung}

\skillcard
{\term{Meditation zur Heilung}{meditation-heilung}}
{Bonusaktion}
{4 AP}
{Selbst}
{—}
{
	\begin{itemize}
		\item Stellt \textbf{3 + Rang LP} wieder her.
		\item Nur auf den Mönch selbst anwendbar.
		\item Nur \textbf{einmal alle 3 Runden}.
	\end{itemize}
}
{
	Kann als \textbf{Aktion} genutzt werden, um zusätzlich
	bis zu \textbf{2 Verbündete} zu heilen.
}

\paragraph*{Meditation des Schutzes}

\spellcard
{\term{Meditation des Schutzes}{meditation-schutz}}
{Bonusaktion}
{4 MP}
{Selbst}
{Bis Kampfende}
{
	\begin{itemize}
		\item Erzeugt \textbf{Magisches Schild} in Höhe von \textbf{2 × Rang}.
		\item Nur \textbf{einmal alle 3 Runden}.
		\item Am Kampfende wird das Schild entfernt.
	\end{itemize}
}
{
	Nahe Verbündete (1 m) profitieren vom Schild,
	erhalten jedoch \textbf{kein eigenes SP}.
}

\paragraph*{Reinigende Hand}

\spellcard
{\term{Reinigende Hand}{reinigende-hand}}
{Aktion}
{4 MP}
{Berührung}
{—}
{
	\begin{itemize}
		\item Wähle \textbf{Statuseffekt oder Zusatzeffekt}.
		\item Verursacht \textbf{2 + Anzahl vorhandener Effekte Direktschaden}.
		\item Entfernt gleichzeitig bis zu \textbf{Rang Effekte}.
	\end{itemize}
}
{
	Entfernte Status- und Zusatzeffekte können
	für den restlichen Kampf \textbf{nicht erneut} angewendet werden.
}

\paragraph*{Aura der Ruhe}

\spellcard
{\term{Aura der Ruhe}{aura-der-ruhe}}
{Bonusaktion}
{$5 - \text{Rang}$ MP}
{Radius 3 m}
{2 Runden}
{
	\textbf{Wähle einen Effekt:}
	\begin{itemize}
		\item \textbf{Gegner:} --1 auf Angriffs- und Verteidigungsproben
		(KAMPF-Probe negiert Effekt).
		\item \textbf{Verbündete:} +1 auf Angriffs- und Verteidigungsproben.
	\end{itemize}
}
{
	Gegner verlieren \textbf{5 AP/MP}
	\textbf{oder}
	Verbündete erhalten \textbf{5 AP/MP}.
}

\paragraph*{Erleuchteter Schlag}

\spellcard
{\term{Erleuchteter Schlag}{erleuchteter-schlag}}
{Aktion}
{8 MP}
{Nahkampf}
{—}
{
	\begin{itemize}
		\item Verursacht \textbf{2 × Rang + STÄRKE-Wert Schaden}.
		\item Tötet der Angriff das Ziel:
		\textbf{Heilung 1W6 LP} für den Mönch.
	\end{itemize}
}
{
	Gib \textbf{1 Marke} aus:
	Der Schaden wird als \textbf{Direktschaden} behandelt.
}

\subsection{Alchemist}

\subsubsection*{Rolle}
Unterstützer / Taktiker

\subsubsection*{Kurzprofil}
Der \term{Alchemist}{klasse-alchemist} ist ein taktischer Unterstützer,
der Tränke, Gifte, Fallen und improvisierte Konstruktionen nutzt,
um Verbündete zu stärken und Gegner gezielt zu kontrollieren.

\subsubsection*{Stärken}
\begin{itemize}
	\item Heilung und Buffs über Tränke
	\item Gifte, Explosionen und Fallen
	\item Hohe Flexibilität durch Materialien
	\item Starke Kontrolle ohne direkte Frontlinie
\end{itemize}

\subsubsection*{Kampfstil}
Indirekt, taktisch, ressourcenbasiert.

% =========================
% Startausrüstung
% =========================
\subsubsection*{Startausrüstung}

\begin{center}
	\begin{tabular}{lp{9cm}}
		\toprule
		\textbf{Gegenstand} & \textbf{Wert und Eigenschaften} \\
		\midrule
		Leichte Rüstung (normal) & 2 Rüstungspunkte (RP) \\
		Hammer & STÄRKE +1 Schaden \\
		Rundschild (normal) & 2 Rüstungspunkte (RP) \\
		Heiltrank & Siehe Tränke \\
		Gifttrank oder Explosiver Trank & Spielerwahl \\
		Alchemieausrüstung & \textbf{10 Materialien} \\
		\bottomrule
	\end{tabular}
\end{center}

Der Alchemist startet jeden Tag mit 10 Materialien, die für die Ausführung seiner Skills
erforderlich sind.

\begin{rulebox}{Materialien}
	Materialien sind eine spezielle Ressource für den Alchemisten.
	Sie werden für Tränke, Fallen und Verbesserungen verbraucht.
	Details siehe \textbf{Ressourcen-Kapitel}.
\end{rulebox}

% =========================
% Klassenfähigkeiten
% =========================
\subsubsection*{Klassenfähigkeiten}

\paragraph*{Giftiger Messerwurf}

\skillcard
{\term{Giftiger Messerwurf}{giftiger-messerwurf}}
{Bonusaktion}
{4 AP}
{Wurfwaffe}
{—}
{
	\begin{itemize}
		\item Verursacht \textbf{Rang Schaden}.
		\item Ziel erleidet \textbf{Vergiftung Rang 2} für \textbf{2 Runden}.
	\end{itemize}
}
{
	Vergiftung steigt auf \textbf{Rang 5} für \textbf{3 Runden}.
}

\paragraph*{Falle platzieren}

\skillcard
{\term{Falle platzieren}{falle-platzieren}}
{Aktion}
{4 AP + 2 Materialien}
{Nahbereich}
{Bis Auslösung}
{
	\begin{itemize}
		\item Platzierung dauert \textbf{2 Runden}.
		\item Gegner müssen eine \textbf{GEWANDTHEIT-Probe – Rang} bestehen.
		\item Die Falle bleibt bis Auslösung oder Kampfende aktiv.
	\end{itemize}
	
	\textbf{Mögliche Fallenarten:}
	\begin{itemize}
		\item Betäubungsfalle – Betäubung (2 Runden)
		\item Explosionsfalle – 1W8 Schaden
		\item Ablenkungsfalle – Verwirrung (2 Runden)
		\item Lichtfalle – Blendung Rang
		\item Eisfalle – Unterkühlung Rang
		\item Feuerfalle – Verbrennung Rang
	\end{itemize}
}
{
	Es können \textbf{zwei Fallen} zu einer kombiniert werden.
}

\paragraph*{Alchemistische Verbesserung}

\skillcard
{\term{Alchemistische Verbesserung}{alchemistische-verbesserung}}
{Aktion}
{Siehe Tabelle}
{Berührung}
{Siehe Tabelle}
{
	\begin{center}
		\begin{tabular}{cccc}
			\toprule
			\textbf{Rang} & \textbf{Verbesserung} & \textbf{Dauer} & \textbf{Kosten} \\
			\midrule
			1 & +1 Schaden \textbf{oder} +1 RP & 1 Runde & 5 AP \\
			2 & +2 Schaden \textbf{oder} +2 RP & 1 Runde & 5 AP \\
			3 & +2 Schaden \textbf{und} +2 RP & 2 Runden & 5 AP \\
			4 & +2/+2 + dauerhaft +1 & 2 Runden & 8 AP + 1 Mat. \\
			5 & +2/+2 + dauerhaft +2 & 2 Runden & 8 AP + 1 Mat. \\
			\bottomrule
		\end{tabular}
	\end{center}
}
{
	Dauerhafte Verbesserungen können einmalig auf einen
	\textbf{Elementarschaden} gepolt werden
	(Rüstung erhält Resistenz).
}

\paragraph*{Materialien sammeln}

\skillcard
{\term{Materialien sammeln}{materialien-sammeln}}
{Aktion}
{2 AP}
{Umgebung}
{—}
{
	\begin{itemize}
		\item Sammelt \textbf{2 × Rang Materialien}.
	\end{itemize}
}
{
	Gegenstände können zerlegt und in Materialien umgewandelt werden
	(Menge nach SL-Entscheid).
}

\paragraph*{Tränke herstellen}

\skillcard
{\term{Tränke herstellen}{traenke-herstellen}}
{Aktion}
{3 Aktionspunkte (AP) + 1 Material}
{Berührung}
{Sofort (Stufe 1) / 2 Runden (Stufe 2)}
{
	\begin{itemize}
		\item Stelle einen Trank her und gib ihn dir selbst oder einem Verbündeten.
		\item \textbf{Grundkosten:} 3 AP und \textbf{1 Material}.
		\item \textbf{Stufe 1:} Der Trank ist \textbf{sofort} hergestellt (in dieser Aktion).
		\item \textbf{Stufe 2:}
		\begin{itemize}
			\item kostet \textbf{zusätzliche Ressourcen} (siehe Trankliste),
			\item benötigt \textbf{eine weitere Runde} (insgesamt 2 Runden Herstellzeit).
		\end{itemize}
		\item Welche Tränke verfügbar sind, richtet sich nach dem \textbf{Rang} dieses Alchemisten-Skills
		(siehe Tranklisten nach Rang).
	\end{itemize}
}
{
	Tränke bis \textbf{Rang 3} können als \textbf{Bonusaktion} hergestellt werden.
}

\medspace

\textbf{Verfügbare Tränke nach Rang:}

\begin{center}
	\begin{tabularx}{\linewidth}{c l X X}
		\toprule
		\textbf{Rang} &
		\textbf{Trank} &
		\textbf{Stufe 1 } &
		\textbf{Stufe 2 } \\
		\midrule
		1 &
		Heiltrank &
		Stellt \textbf{5 LP} wieder her &
		Stellt \textbf{10 LP} wieder her \\
		\midrule
		2 &
		Manatrank &
		Stellt \textbf{10 MP} wieder her &
		Stellt \textbf{15 MP} wieder her \\
		2 &
		AP-Trank &
		Stellt \textbf{10 AP} wieder her &
		Stellt \textbf{15 AP} wieder her \\
		\midrule
		3 &
		Gifttrank &
		\textbf{Vergiftung Rang 1} + \textbf{1W4 Schaden} (Werfen) &
		\textbf{Vergiftung Rang 2} + \textbf{1W6 Schaden} (Werfen) \\
		3 &
		Explosiver Trank &
		\textbf{1W4 Feuerschaden} (2\,m Radius) &
		\textbf{1W6 Feuerschaden} (2\,m Radius) \\
		\midrule
		4 &
		Kombitrank &
		Stellt \textbf{je 3 LP, MP und AP} wieder her &
		Stellt \textbf{je 6 LP, MP und AP} wieder her \\
		4 &
		Stärketrank &
		\textbf{+2 Bonus auf STÄRKE} für 3 Runden &
		— \\
		\midrule
		5 &
		Rasereielixier &
		\textbf{+2 Bonus auf Nahkampf} und \textbf{+2 Bonus auf STÄRKE} für 3 Runden &
		— \\
		5 &
		Antimagisches Serum &
		\textbf{–5 Malus auf alle Zauber} für 10 Minuten &
		— \\
		\bottomrule
	\end{tabularx}
\end{center}

\paragraph*{(fehlender 6. Zauber)}
% TODO:  6. Klassenzauber liefern

\section{Klassen (magisch)}
\subsection{Magische Basisausrüstung}

\term{Magische Klassen}{magierklassen} erhalten \textbf{keine fest zugewiesene Startausrüstung}.
Stattdessen wählen sie bei der Charaktererschaffung
\textbf{eine der folgenden magischen Basisausrüstungen}.

Diese Wahl ist unabhängig von der Klasse,
sofern eine Klasse keine explizite Einschränkung angibt.

% =========================
% Magischer Gelehrter
% =========================
\subsubsection*{\term{Der magische Gelehrte}{ausruestung-magischer-gelehrter}}

\begin{rulebox}{Basisausrüstung – Magischer Gelehrter}
	\begin{tabular}{lp{9cm}}
		\toprule
		\textbf{Gegenstand} & \textbf{Regelwirkung} \\
		\midrule
		Zauberbuch &
		\textbf{+1 Bonus auf WISSEN-WURFZAHL} \\
		
		Kristallfokus &
		\textbf{1× pro Tag: Magie identifizieren}
		(\textbf{WISSEN}-Probe) \\
		
		Einfacher magischer Dolch &
		\textbf{WISSEN + 1 Schaden} \\
		
		\bottomrule
	\end{tabular}
\end{rulebox}

\textbf{Rolle:} Arkaner Wissenshüter\\  
\textbf{Schwerpunkt:} Analyse, Identifikation, magische Aufklärung\\  
\textbf{Spielweise:} Defensiv, unterstützend, zweite Reihe\\

% =========================
% Kampfmagier
% =========================
\subsubsection*{\term{Der Kampfmagier}{ausruestung-kampfmagier}}

\begin{rulebox}{Basisausrüstung – Kampfmagier}
	\begin{tabular}{lp{9cm}}
		\toprule
		\textbf{Gegenstand} & \textbf{Regelwirkung} \\
		\midrule
		Magische leichte Kampfrüstung (normal) &
		\textbf{2 Rüstungspunkte (RP)},
		\textbf{4 Magisches Schild (SP)} \\
		
		Kurzschwert &
		\textbf{STÄRKE + 2 Schaden} \\
		
		2× Manatrank (Stufe 1) &
		Jeweils \textbf{10 Manapunkte (MP)} \\
		
		Notfall-Schriftrolle &
		\textbf{Explosion:}
		Radius 2 m, \textbf{2W6 Schaden} \\
		
		\bottomrule
	\end{tabular}
\end{rulebox}

\textbf{Rolle:} Frontzauberer / Arkaner Krieger\\  
\textbf{Schwerpunkt:} Direkter Kampf, Schadenszauber, Durchhaltefähigkeit\\  
\textbf{Spielweise:} Aggressiv, flexibel, vordere Linie

% =========================
% Magischer Unterstützer
% =========================
\subsubsection*{\term{Der magische Unterstützer}{ausruestung-magischer-unterstuetzer}}

\begin{rulebox}{Basisausrüstung – Magischer Unterstützer}
	\begin{tabular}{lp{9cm}}
		\toprule
		\textbf{Gegenstand} & \textbf{Regelwirkung} \\
		\midrule
		Einfacher magischer Dolch &
		\textbf{WISSEN + 1 Schaden} \\
		
		Schwere magische Robe (normal) &
		\textbf{Nicht-Schadenszauber kosten 1 MP weniger}, \\
		& \textbf{6 Magisches Schild (SP)} \\
		
		Anhänger der Geisterfauna &
		\textbf{1× pro Tag:}
		Kommunikation mit einer wilden Kreatur \\
		
		\bottomrule
	\end{tabular}
\end{rulebox}

\textbf{Rolle:} Spiritueller Verbündeter / Arkaner Unterstützer\\  
\textbf{Schwerpunkt:} Buffs, Schutz, Kommunikation, Umweltinteraktion\\  
\textbf{Spielweise:} Zurückhaltend, unterstützend, kontrollierend

\subsection{Elementarmagie}

\subsubsection*{Rolle}
Magischer Manipulator der Naturgewalten

\subsubsection*{Kurzprofil}
Der \term{Elementarmagier}{klasse-elementarmagie} ist ein Meister der rohen Urkräfte.
Er lenkt Feuer, Eis, Erde, Blitz und Wind gleichermaßen, um das Schlachtfeld zu formen,
Gegner zu kontrollieren oder Verbündete zu unterstützen.

\subsubsection*{Stärken}
\begin{itemize}
	\item Sehr vielseitige Zauberauswahl
	\item Starke Flächenkontrolle
	\item Hoher Schaden und situative Defensive
	\item Zugriff auf elementare Resistenz
\end{itemize}

\subsubsection*{Kampfstil}
Flexibel und taktisch – wechselt zwischen mächtiger Offensive
(z.\,B. Elektrische Entladung) und kontrollierenden oder defensiven Effekten
(z.\,B. Steinwall erschaffen).

\subsubsection*{Magische Basisausrüstung (Empfehlung)}
Empfohlene Startausrüstung:
\termref[Der Kampfmagier]{ausruestung-kampfmagier}

% =========================
% Klassenzauber
% =========================
\subsubsection*{Klassenzauber}

\paragraph*{Flammenheilung}
\spellcard
{\term{Flammenheilung}{flammenheilung}}
{Aktion}
{4 MP + 1 MP pro weiterem Ziel}
{5 m}
{Sofort}
{
	Heilt Ziele durch kontrollierte, reinigende Flammenenergie.
	
	\begin{center}
		\begin{tabular}{p{1cm}p{7cm}p{5cm}}
			\toprule
			\textbf{Rang} & \textbf{Heilung} & \textbf{Risiko bei Misserfolg} \\
			\midrule
			1 & 1W4 LP (nur Selbst) & — \\
			2 & 1 + 2W4 LP (bis 2 Ziele) & 1W4 Feuerschaden \\
			3 & 2 + 3W4 LP (bis 2 Ziele) & 2W4 Feuerschaden \\
			4 & 3 + 3W4 LP (bis 3 Ziele) & 1W4 + 1W6 Feuerschaden \\
			5 & 3 + 3W6 LP (bis 3 Ziele) & 1W4 + 2W6 Feuerschaden \\
			\bottomrule
		\end{tabular}
	\end{center}
	
	Zusätzlich erhält der Anwender \textbf{Resistenz gegen Feuerschaden}
	bis Kampfende oder 10 Minuten.
}
{
	\begin{itemize}
		\item \textbf{Phoenixfunke:} Alle geheilten Ziele können in diesem Kampf
		einmalig nicht unter \textbf{1 LP} fallen.
		\item \textbf{Hitzeheilung:} Unterkühlung und Vergiftung werden entfernt.
	\end{itemize}
}

\paragraph*{Froststrahl}
\spellcard
{\term{Froststrahl}{froststrahl}}
{Aktion}
{4 MP}
{15 m}
{Sofort}
{
	Verursacht \textbf{Wurfzahl Kälteschaden Rang 2}.
}
{
	Verursacht zusätzlich \textbf{Erosion Rang 3}.
}

\paragraph*{Steinwall erschaffen}
\spellcard
{\term{Steinwall erschaffen}{steinwall}}
{Aktion}
{8 MP}
{20 m}
{1 Runde pro Rang}
{
	\begin{itemize}
		\item Erschafft einen \textbf{Rang × Rang Meter} großen Steinwall.
		\item Lebenspunkte des Walls: \textbf{10 + 4 × Rang}.
		\item Blockiert Bewegung und Sicht.
	\end{itemize}
}
{
	\begin{itemize}
		\item Der Steinwall ist \textbf{resistent gegen \termref[Elementarschaden]{elementarschaden}}.
		\item Die Lebenspunkte des Walls werden auf \textbf{70 LP} erhöht.
	\end{itemize}
}

\paragraph*{Elektrische Entladung}
\spellcard
{\term{Elektrische Entladung}{elektrische-entladung}}
{Aktion}
{8 MP}
{20 m}
{Sofort}
{
	\begin{itemize}
		\item Verursacht \textbf{WISSEN-Wert + Wurfzahl} Blitzschaden.
		\item Ab \textbf{Wurfzahl 8}: Kettenblitz
		\begin{itemize}
			\item 1W6 Chance, alle weiteren Ziele im Umkreis von
			\textbf{Rang Metern} zu treffen.
			\item Diese erleiden \textbf{50\,\% des Schadens}.
		\end{itemize}
	\end{itemize}
}
{
	\textbf{Kettenblitzgarantie:}
	Jeder Treffer erzeugt automatisch einen Kettenblitz,
	der allen getroffenen Zielen \textbf{vollen Schaden} zufügt.
}

\paragraph*{Windstoß}
\spellcard
{\term{Windstoß}{windstoss}}
{Bonusaktion}
{4 MP}
{10 m}
{Sofort}
{
	\begin{itemize}
		\item Erzeugt eine Windböe im Umkreis von \textbf{5 Metern}.
		\item Ziele werden um \textbf{2 × Rang Meter} zurückgestoßen.
		\item Ab Rang 3: je \textbf{1W6 Chance} ein Ziel zu entwaffnen. Für jedes Ziel muss einzeln gewürfelt werden.
	\end{itemize}
}
{
	\textbf{Sogstoß:}
	Erzeugt einen Sog im Umkreis von \textbf{10 Metern}.
	Alle Ziele im Radius von \textbf{6 Metern} werden
	zum Zentrum gezogen und entwaffnet alle Ziele . 
}

\paragraph*{Blendlicht}
\spellcard
{\term{Blendlicht}{blendlicht}}
{Bonusaktion}
{2 MP}
{10 m}
{1 Runde}
{
	\begin{itemize}
		\item Ziel erhält \termref[Blendung]{blendung} des Rangs.
	\end{itemize}
}
{
	Ziel erhält zusätzlich \termref[Verwirrung]{verwirrung} und wird entwaffnet.
}

\subsection{Tiermagie}

\subsubsection*{Rolle}
Verbündeter der Natur und ihrer Kreaturen

\subsubsection*{Kurzprofil}
Die \term{Tiermagie}{klasse-tiermagie} verbindet den Zaubernden direkt mit der wilden Welt:
Tiere rufen, Befehle geben, Tiergestalten annehmen und Pflanzen beschleunigt wachsen lassen.
Im Kampf ist sie unterstützend oder kontrollierend, in der Wildnis extrem stark.

\subsubsection*{Stärken}
\begin{itemize}
	\item Unterstützung durch Tiere (Aufklärung, Kontrolle, Druck)
	\item Körperliche Verstärkung durch Verwandlung
	\item Starke Wildnis-/Erkundungsoptionen
\end{itemize}

\subsubsection*{Kampfstil}
Unterstützend oder wild -- zwischen Rudeltaktik, Kontrolle und Tiergestalt.

\subsubsection*{Magische Basisausrüstung (Empfehlung)}
Empfohlen: \termref[Der magische Unterstützer]{ausruestung-magischer-unterstuetzer}.  
(Alternativ je nach Spielstil: \termref[Der magische Gelehrte]{ausruestung-magischer-gelehrter}.)

\subsubsection*{Klassenzauber}

\paragraph*{Tierbefehl}
\spellcard
{\term{Tierbefehl}{tierbefehl}}
{Aktion}
{4 MP}
{5 Kilometer}
{Bis erfüllt, max. Wurfzahl Stunden}
{
	\begin{itemize}
		\item Es erscheint ein Tier der gewünschten Spezies (nur wenn Tiere dieser Spezies in der Nähe existieren).
		\item Das Tier führt \textbf{einen Befehl} aus (\textbf{kein Angriff}).
		\item Die \textbf{Wurfzahl} bestimmt die Effektivität; Schwellen für außergewöhnliche Aufgaben legt die SL fest.
		\item Kurze Befehle können auch als \textbf{Bonusaktion} erteilt werden.
	\end{itemize}
}
{
	\textbf{Lautlose Befehle:} Du kannst nun bis zu \textbf{2 Befehle gleichzeitig} \textbf{lautlos} erteilen.
}

\paragraph*{Tiersprache}
\spellcard
{\term{Tiersprache}{tiersprache}}
{Bonusaktion}
{2 MP}
{1/5/25/100/1000 Meter}
{Bis das Tier das Interesse verliert}
{
	\begin{itemize}
		\item Du kannst mit einem Tier sprechen.
		\item Je höher die \textbf{Wurfzahl}, desto hilfreicher/konkreter sind die Informationen.
		\item Besonderheit Reichweite: Abhängig vom Rang kann ein entferntes Tier \textbf{telepathisch gerufen} werden (siehe Reichweite).
		\item Je höher der \textbf{Rang}, desto länger bleibt das Tier am Gespräch interessiert.
	\end{itemize}
}
{
	\textbf{Empathischer Blick:} Du erkennst bei einem Tier in \textbf{5 m} sofort Emotion und Grundzustand
	(\textit{verletzt/hungrig/ängstlich}).
}

\paragraph*{Tierverwandlung}
\spellcard
{\term{Tierverwandlung}{tierverwandlung}}
{Aktion}
{16 - Rang MP}
{Selbst}
{Bis beendet, spätestens nächster Tag}
{
	\begin{itemize}
		\item Du verwandelst dich in ein Tier und übernimmst dessen Fähigkeiten.
		\item In Tiergestalt kannst du \textbf{nicht zaubern}.
		\item Du erhältst \textbf{separate LP}; beim Ende der Verwandlung werden deine LP wieder auf deine normalen LP zurückgesetzt.
		\item Kleidung/Ausrüstung verwandeln sich mit.
		\item Für Beispiele von Tiere nach Rang siehe Tabelle unten.
	\end{itemize}
	
}
{
	\textbf{Hybridform:} Du kannst dich in eine Hybridgestalt aus \textbf{zwei Tieren} verwandeln.
	Als Werte können die \textbf{Maximalwerte} der jeweiligen Würfel genommen werden (oder normal gewürfelt).
}

\textbf{Tierformen nach Rang:}

\medskip

\textbf{Rang 1 und 2}
\creatureblock
{Ameise}
{1 LP / 0 RP / 0 SP}
{STÄRKE 1W4, KAMPF 1W4, GEWANDTHEIT 1W12}
{Biss: Schaden = STÄRKE-Wert / 2}
{Winzig: Kann durch kleinste Öffnungen gelangen}

\creatureblock
{Maus}
{3 LP / 0 RP / 0 SP}
{STÄRKE 1W4, KAMPF 1W4, GEWANDTHEIT 1W10}
{Biss: Schaden = STÄRKE-Wert / 2}
{Unauffällig: Vorteil auf Heimlichkeitsproben}

\textbf{Rang 3 und 4}
\creatureblock
{Vogel}
{5 LP / 0 RP / 0 SP}
{STÄRKE 1W6, KAMPF 1W6, GEWANDTHEIT 1W8}
{Schnabel: Schaden = STÄRKE-Wert}
{Flugfähig}

\creatureblock
{Hund}
{15 LP / 0 RP / 0 SP}
{STÄRKE 1W8, KAMPF 1W8, GEWANDTHEIT 1W6}
{Biss: Schaden = STÄRKE-Wert}
{Spürnase: Vorteil auf Wahrnehmung (Geruch)}

\textbf{Rang 5}
\creatureblock
{Bär}
{40 LP / 5 RP / 0 SP}
{STÄRKE 1W10, KAMPF 1W8, GEWANDTHEIT 1W4}
{Pranke: Schaden = 2× STÄRKE-Wert}
{Massiv: Immun gegen Niederschlag}


\paragraph*{Pflanzenwachstum}
\spellcard
{\term{Pflanzenwachstum}{pflanzenwachstum}}
{Aktion}
{4 MP}
{10 Meter}
{1 Minute pro Rang}
{
	\begin{itemize}
		\item Pflanzen wachsen innerhalb weniger Sekunden extrem schnell.
		\item Erzeugt \textbf{dichte Sträucher/Bäume/Ranken} als Hindernis oder Deckung.
		\item Größe: Höhe und Breite = \textbf{Rang Meter}.
	\end{itemize}
}
{
	\textbf{Giftwuchs:} Die Pflanzen sind giftig und verursachen \termref[Vergiftung]{vergiftung} Rang 5 für \textbf{3 Runden}.
}

\paragraph*{Schwarm beschwören}
\spellcard
{\term{Schwarm beschwören}{schwarm-beschwoeren}}
{Aktion}
{8 MP}
{20 Meter}
{1 Runde}
{
	\begin{itemize}
		\item Du beschwörst einen Schwarm aus Insekten oder kleinen Tieren.
		\item Der Schwarm attackiert ein Ziel und verursacht \textbf{Rang + Wurfzahl Schaden}.
	\end{itemize}
}
{
	Das Ziel erhält zusätzlich \termref[Blendung]{blendung}, \termref[Verwirrung]{verwirrung} und \termref[Erschöpfung]{erschoepfung}.
}

\paragraph*{Ruf des Tiers}
\spellcard
{\term{Ruf des Tiers}{ruf-des-tiers}}
{Aktion}
{8 MP}
{3 Meilen}
{2 + Rang Runden}
{
	\begin{itemize}
		\item Ein wildes Tier erscheint und unterstützt dich im Kampf.
		\item Die \textbf{Wurfzahl} beeinflusst, ob ein stärkeres oder schwächeres Tier erscheint (SL-Entscheid).
		\item Je höher der \textbf{Rang}, desto loyaler ist das Tier; schwächere Tiere fliehen bei hohem Schaden.
	\end{itemize}
}
{
	Zusätzlich Zugriff auf legendäre Tiere unabhängig vom Aufenthaltsort (siehe Tabelle unten)
}

\textbf{Legendäre Tiere (Rang 6):}

\medskip

\creatureblock
{Schattenwolf – Jäger aus der Finsternis}
{18 LP / 2 RP / 0 SP}
{KAMPF 5, GEWANDTHEIT 4, WAHRNEHMUNG 4}
{%
	Biss: 1W8 Schaden \newline
	Jagdheulen (3, Aktion): Ziel markieren → +1 Angriff für Verbündete \newline
	Giftbiss (4, Aktion, 1×/Ziel/Kampf): 1W6 Giftschaden, Vergiftung Rang 3 (3 Runden)
}
{Schattenhaut: Nach Treffer 1 m Bewegung ohne Gelegenheitsangriffe}

\medskip

\creatureblock
{Himmelsfalke – Auge von oben}
{12 LP / 1 RP / 0 SP}
{KAMPF 4, GEWANDTHEIT 5, WAHRNEHMUNG 5}
{%
	Klauenstoß: 1W6 Schaden \newline
	Markieren (2, Bonus): Nächster Fernangriff ignoriert SP \newline
	Sturzflug (4, Aktion): Blendung Rang 3 oder Verlust Bonusaktion
}
{Hochauge: 1 klarer Hinweis pro Szene durch Luftaufklärung}

\medskip

\creatureblock
{Schwarzbär – Lebende Bastion}
{30 LP / 5 RP / 0 SP}
{KAMPF 5, STÄRKE 5, WAHRNEHMUNG 3}
{%
	Pranke: 1W10 Durchschlag \newline
	Prankenstoß (4, Aktion): +1W4 Durchschlag + Niederschlag oder 1 m Schieben \newline
	Wächterfell (3, Bonus): Fängt Treffer für Verbündeten in 1 m ab
}
{Panzerfell: Kein Durchschlagsschaden möglich}

\subsection{Zauberschmied}

\subsubsection*{Rolle}
Runenmeister und Magiehandwerker

\subsubsection*{Kurzprofil}
Der \term{Zauberschmied}{klasse-zauberschmied} verbindet arkane Macht mit handwerklicher Meisterschaft.
Er erschafft, verzaubert und verändert Waffen, Rüstungen und magische Gegenstände.
Seine Macht liegt weniger in der flüchtigen Geste, sondern im geschmiedeten Werk.

\subsubsection*{Stärken}
\begin{itemize}
	\item Anpassbare Ausrüstung und Verzauberungen
	\item Starke Vorbereitungsmöglichkeiten
	\item Unterstützung für Verbündete durch Gegenstands-Magie
	\item Kann auch im Nahkampf bestehen (über Ausrüstung)
\end{itemize}

\subsubsection*{Kampfstil}
Reaktiv und unterstützend – optimiert Waffen und Rüstung für die jeweilige Situation.

\subsubsection*{Magische Basisausrüstung (Empfehlung)}
Empfohlen: \termref[Der magische Unterstützer]{ausruestung-magischer-unterstuetzer}.  
(Alternativ je nach Spielstil: \termref[Der magische Gelehrte]{ausruestung-magischer-gelehrter}.)

\subsubsection*{Klassenzauber}

\paragraph*{Magische Reparatur}
\spellcard
{\term{Magische Reparatur}{zauber-magische-reparatur}}
{Aktion}
{5 MP}
{5 Meter}
{—}
{
	\begin{itemize}
		\item Repariert ein beschädigtes oder zerstörtes \textbf{nichtmagisches} Objekt.
		\item Kann auch \textbf{ \termref[Rüstung]{ruestung} (RP)} wiederherstellen.
		\item Je größer der Schaden, desto höher die erforderliche Schwelle für die \textbf{Wurfzahl} (SL-Entscheid).
		\item \textbf{Magische Gegenstände} können erst ab \textbf{Rang 3} repariert werden.
	\end{itemize}
	
	\medskip
	\textbf{Zauberdauer:} 5 Minuten
}
{
	\begin{itemize}
		\item Gegenstände werden \textbf{sofort vollständig} wiederhergestellt.
		\item Gib \textbf{1 Marke} aus, um die \textbf{Qualität einer Rüstung (RP)} für den nächsten Kampf zu erhöhen.
		\item Dauerhafte Erhöhung kostet \textbf{doppelte Manapunkte (MP)}.
	\end{itemize}
}

\paragraph*{Verzauberung}
\spellcard
{\term{Verzauberung}{zauber-verzauberung}}
{Aktion}
{8 MP}
{Selbst (Objekt berühren)}
{Langfristig}
{
	\begin{itemize}
		\item Verzaubert einen Gegenstand, um \termref[Schadenstypen]{schadenstypen} oder \termref[Statuseffekte]{statuseffekte} zu erzeugen.
		\item Ab \textbf{Rang 4} sind komplexere Verzauberungen möglich, z.\,B.:
		\begin{itemize}
			\item Kombination mehrerer Schadenstypen (z.\,B. Feuer + Blitz).
			\item Verzauberungen mit \textbf{auslösenden Bedingungen}
			(z.\,B. Effekt triggert nur bei Treffer, Parade oder kritischem Erfolg).
			\item Verzauberungen mit \textbf{begrenzten Ladungen} oder Abklingzeiten.
			\item Defensive Effekte mit Rückwirkung (z.\,B. Schaden bei Treffer auf den Angreifer).
		\end{itemize}
		\item Bestehende Verzauberungen können \textbf{nicht} überschrieben werden.
	\end{itemize}
	
	\medskip
	\textbf{Zauberdauer:} 1--10 Stunden (je nach Rang und Komplexität)
}
{
	\begin{itemize}
		\item Einfache Änderungen (Schadenstyp oder Zusatzeffekt) können als \textbf{Bonusaktion} und ohne Zauberdauer ausgeführt werden.
		\item Bestehende Verzauberungen können nun \textbf{überschrieben} werden.
	\end{itemize}
}

\paragraph*{Zauberspeicher}
\spellcard
{\term{Zauberspeicher}{zauber-zauberspeicher}}
{Aktion (Speichern), Bonusaktion (Aktivieren)}
{6 MP beim Speichern; Aktivieren kostenlos}
{Selbst (Objekt berühren)}
{Bis Auslösung}
{
	\begin{itemize}
		\item Lädt einen Gegenstand mit einem Zauber, der später ausgelöst werden kann.
		\item Der Rang des gespeicherten Zaubers ist \textbf{immer 1 niedriger} als der Originalzauber.
		\item Es können nur Zauber gespeichert werden, die \textbf{gleich oder niedriger} sind als der Rang von Zauberspeicher.
		\item Wird der Gegenstand zerstört, werden gespeicherte Zauber \textbf{automatisch ausgelöst}.
		\item Schaden, der durch \textbf{Wurfzahl} definiert wird, wird nicht gewürfelt:
		es gilt der \textbf{Durchschnittswert} (z.\,B. bei Rang 4/1W10 ist es automatisch 5).
	\end{itemize}
}
{
	\begin{itemize}
		\item Zauber können nun im \textbf{Originalrang} gespeichert werden.
		\item \textbf{Verdopplung:} Gespeicherte Zauber von Rang 4 oder niedriger können
		\textbf{zweimal aktiviert} werden oder bei einer Aktivierung \textbf{verdoppelt} werden.
	\end{itemize}
}

\paragraph*{Energieschild}
\spellcard
{\term{Energieschild}{zauber-energieschild}}
{Aktion}
{2 MP pro Rang}
{10 Meter}
{Bis Kampfende oder 1 Minute pro Rang}
{
	\begin{itemize}
		\item Erzeugt ein \termref[Magisches Schild (SP)]{magischer-schild} auf dem Ziel.
		\item höhere Ränge können bestimmte Schadenstypen vollständig blockieren (keinen Schaden)
		\item Effekte skalieren nach Rang:
	\end{itemize}
	
	\begin{center}
		\begin{tabular}{lp{10cm}}
			\toprule
			\textbf{Rang} & \textbf{Effekt} \\
			\midrule
			1 & 2 Magisches Schild (SP) \\
			2 & 3 Magisches Schild (SP), blockiert zusätzlich \termref[Feuerschaden]{feuerschaden} \\
			3 & 4 Magisches Schild (SP), blockiert zusätzlich  \termref[Elementarschaden]{elementarschaden} \\
			4 & 6 Magisches Schild (SP), blockiert zusätzlich \termref[Elementarschaden]{elementarschaden};
			\newline einmalige 1W4-Chance, dass der Schild nur auf \textbf{1 LP} fällt \\
			5 & 8 Magisches Schild (SP), blockiert zusätzlich \termref[Elementarschaden]{elementarschaden};
			\newline kann einmalig durch einen einzelnen Angriff nur auf \textbf{1 LP} fallen \\
			\bottomrule
		\end{tabular}
	\end{center}
}
{
	\textbf{Spiegelschild:} Gib \textbf{2 Marken} aus, um jeglichen eingehenden Schaden auf das Energieschild
	in dieser Runde auf den Ursprung zurückzuleiten.
}

\paragraph*{Arkaner Spürsinn}
\spellcard
{\term{Arkaner Spürsinn}{zauber-arkaner-spuersinn}}
{Bonusaktion}
{$6 - \text{Rang}$ Manapunkte (MP)}
{30 Meter}
{1 Minute}
{
	\begin{itemize}
		\item Erkennt magische Signaturen im Umkreis: \textbf{Gegenstände}, \textbf{Fallen} und \textbf{aktive Zauber}.
		\item Kann \textbf{spezifische Energien} und die \textbf{Stärke} der Magie unterscheiden.
	\end{itemize}
}
{
	Gib \textbf{1 Marke} aus, um eine automatische Erkennung aller Magie im Umkreis von \textbf{10 Metern} zu erhalten.
}

\paragraph*{(fehlender 6. Zauber)}
% TODO:  6. Klassenzauber liefern

\subsection{Golemmagie}

\subsubsection*{Rolle}
Beschwörer konstruktiver Wächter

\subsubsection*{Kurzprofil}
Der \term{Golemmagier}{klasse-golemmagie} erschafft und kontrolliert künstliche Wesen aus Stein, Erde oder anderen Materialien.
Seine Macht liegt weniger in direkter Zerstörung, sondern im \textbf{strategischen Einsatz} langlebiger Konstrukte.

\subsubsection*{Stärken}
\begin{itemize}
	\item Hohe Kontrolle über Diener
	\item Starke defensive Möglichkeiten
	\item Gute Unterstützung durch Konstrukte
\end{itemize}

\subsubsection*{Kampfstil}
Strategisch und indirekt – agiert über Golems, positioniert sie taktisch, bleibt selbst meist im Hintergrund.

\subsubsection*{Magische Basisausrüstung (Empfehlung)}
Empfohlen: \termref[Der magische Unterstützer]{ausruestung-magischer-unterstuetzer}.  
(Alternativ je nach Spielstil: \termref[Der Kampfmagier]{ausruestung-kampfmagier}.)

\subsubsection*{Klassenzauber}

\paragraph*{Golem rufen}

\spellcard
{\term{Golem rufen}{golem-rufen}}
{Aktion}
{6 MP (+3 MP klein / +6 MP normal je 30 Minuten)}
{10 m}
{30 Minuten pro Rang}
{
	\begin{itemize}
		\item Beschwört einen Golem, der einfache Befehle ausführt.
		\item Ab \textbf{Rang 2} kann ein normalgroßer Golem erschaffen werden.
		\item Verlängerung kostet erneut die angegebenen Manapunkte.
	\end{itemize}
	
	\medskip
	\textbf{Kleiner Golem:}
	\creatureblock
	{Kleiner Golem}
	{5 LP / 0 RP / 0 SP}
	{STÄRKE 2, KAMPF 2, GEWANDTHEIT 4}
	{Schlag: 1W4 Schaden}
	{--}
	
	\medskip
	\textbf{Normaler Golem:}
	\creatureblock
	{Normaler Golem}
	{15 LP / 0 RP / 0 SP}
	{STÄRKE 3, KAMPF 3, GEWANDTHEIT 3}
	{Schlag: 1W6 Schaden}
	{--}
}
{
	\textbf{Multiplikator:}
	Alle 30 Minuten erzeugt der Golem kostenlos eine exakte Kopie seines aktuellen Zustands
	(Verlängerungskosten bleiben bestehen).
}


\paragraph*{Golemstärke}

\spellcard
{\term{Golemstärke}{golemstaerke}}
{Bonusaktion}
{5 MP}
{10 m}
{1 Runde pro Rang}
{
	\begin{itemize}
		\item Der Golem verursacht \textbf{+Wurfzahl Schaden} mit allen Angriffen.
	\end{itemize}
}
{
	Der gesamte Schaden des Golems wird zu \textbf{Durchschlagsschaden}.
}


\paragraph*{Golemreparatur}

\spellcard
{\term{Golemreparatur}{golemreparatur}}
{Bonusaktion}
{4 MP}
{5 m}
{Sofort}
{
	\begin{itemize}
		\item Heilt \textbf{Wurfzahl / 2} LP (aufgerundet).
		\item Kann keinen vollständig zerstörten Golem reparieren.
		\item RP und SP werden nicht wiederhergestellt.
	\end{itemize}
}
{
	\textbf{Notfallbelebung:}
	Gib \textbf{2 Marken} aus, um einen zerstörten Golem vollständig wiederherzustellen
	(LP, RP und SP).
}


\paragraph*{Wachgolem}

\spellcard
{\term{Wachgolem}{wachgolem}}
{Aktion}
{8 MP}
{10 m}
{2 Runden pro Rang}
{
	\begin{itemize}
		\item Der Wachgolem schützt ein Ziel oder Gebiet (Abwehrmodus).
		\item Verliert jede Runde LP durch \termref[Erosion]{erosion} Rang 1.
	\end{itemize}
	
	\medskip
	\creatureblock
	{Wachgolem}
	{30 LP / Rang RP / Rang SP}
	{STÄRKE 3, KAMPF 4, GEWANDTHEIT 2}
	{Schlag: 1W6 Schaden}
	{Abwehrmodus}
}
{
	\textbf{Abfangmodus:}
	Alle Fernkampfangriffe und Distanzzauber werden automatisch auf den Wachgolem umgeleitet
	und verursachen nur \textbf{50 \% Schaden}.
}


\paragraph*{Körpertransfer}

\spellcard
{\term{Körpertransfer}{koerpertransfer}}
{Aktion}
{10 MP}
{Selbst}
{1 Stunde pro Rang}
{
	\begin{itemize}
		\item Der Zauberer übernimmt vollständig den Körper eines Golems.
		\item Eigene Talente werden durch die des Wirtskörpers ersetzt.
		\item Währenddessen können \textbf{keine Zauber} gewirkt werden.
	\end{itemize}
}
{
	Gib \textbf{1 Marke} aus, um sofort die Position mit dem Ziel zu tauschen.
}


\paragraph*{Erdgolem erschaffen}

\spellcard
{\term{Erdgolem erschaffen}{erdgolem}}
{Aktion}
{12 MP}
{5 m}
{2 Runden pro Rang}
{
	\begin{itemize}
		\item Beschwört einen besonders widerstandsfähigen Erdgolem.
		\item Immun gegen \termref[Durchschlagsschaden]{durchschlagsschaden} (Panzerkörper).
		\item Verliert jede Runde LP durch \termref[Erosion]{erosion} Rang 2.
	\end{itemize}
	
	\medskip
	\creatureblock
	{Erdgolem}
	{25 LP / Rang RP / Rang SP}
	{STÄRKE 5, KAMPF 5, GEWANDTHEIT 3}
	{Pranke: 1W10 Durchschlagsschaden}
	{Panzerkörper}
}
{
	\textbf{Finale:}
	Gib \textbf{2 Marken} aus, um den Kern zu überladen.
	Der Golem explodiert und verursacht Schaden in Höhe seiner maximalen
	\textbf{LP + RP} an allen Zielen im Umkreis von \textbf{15 m}.
}

\subsection{Dunkle Magie}

\subsubsection*{Rolle}
Verbotener Seelenformer / Meister der Schattenkräfte

\subsubsection*{Kurzprofil}
Die \term{Dunkle Magie}{klasse-dunkle-magie} schöpft Macht aus Tod, Schmerz und Schatten.
Ihre Zauber sind effizient, aber gefährlich: Viele Effekte haben direkte Nebenwirkungen
(z.\,B. LP-Verlust, Schadensanfälligkeit, Kontrollrisiko).

\subsubsection*{Stärken}
\begin{itemize}
	\item Sehr hohe Schadenspotenziale und Kontrollzauber
	\item Zugriff auf Untote und Leichenverwertung
	\item Mächtige Trade-offs (Risiko gegen Wirkung)
\end{itemize}

\subsubsection*{Kampfstil}
Berechnend, manipulierend – agiert über Flüche, Kontrolle und Leichenverwertung.

% =========================
% Magische Basisausrüstung (Empfehlung)
% =========================
\subsubsection*{Magische Basisausrüstung (Empfehlung)}
\termref[Der Kampfmagier]{ausruestung-kampfmagier} (Alternativ je nach Spielstil: \termref[Der magische Gelehrte]{ausruestung-magischer-gelehrter}.)

% =========================
% Zusätzliche Ausrüstung (klassen-spezifisch)
% =========================
\subsubsection*{Zusätzliche Ausrüstung}

\begin{center}
	\begin{tabular}{lp{10cm}}
		\toprule
		\textbf{Gegenstand} & \textbf{Eigenschaften} \\
		\midrule
		\term{Ritualdolch}{item-ritualdolch} &
		\textbf{WISSEN +1 Schaden}; kann aus Leichen \textbf{1W4 LP / MP / AP} extrahieren. \\
		\term{Roben der Schatten}{item-roben-der-schatten} &
		\textbf{+2 Bonus auf HEIMLICHKEIT-WURFZAHL}. \\
		\term{Anhänger der Seelen}{item-anhaenger-der-seelen} &
		Tausche \textbf{1 LP} gegen \textbf{2 MP oder 2 AP}. \\
		\bottomrule
	\end{tabular}
\end{center}

% =========================
% Klassenzauber
% =========================
\subsubsection*{Klassenzauber}

\paragraph*{Tote beschwören}

\spellcard
{\term{Tote beschwören}{tote-beschwoeren}}
{Aktion}
{8 MP}
{10 m}
{Bis Tod des Untoten oder Aufhebung (nur bei Kontrolle möglich)}
{
	\begin{itemize}
		\item Beschwört Untote und kontrolliert sie.
		\item \textbf{Kontrollprobe (nur bei Trigger):}
		Sobald du \textbf{Schaden} erleidest oder einen \termref[Statuseffekt]{statuseffekte} erhältst,
		musst du eine neue Probe ablegen (keine erneuten MP-Kosten):
		$$
		Wurfzahl + \lfloor Rang/2 \rfloor \Rightarrow \text{Kontrolle bleibt bestehen}
		$$
		\item \textbf{Kontrollverlust:} Untote greifen willkürlich an (50\% nächstes Ziel,
		ansonsten Verbündete oder den Magier selbst).
		\item Ab \textbf{Rang 2} auch außerhalb des Kampfes nutzbar (Aufgaben/Fragen).
	\end{itemize}
	
	\medskip
	\textbf{Nebenwirkung:}
	\begin{itemize}
		\item Solange der Zauber aktiv ist, sinken deine \textbf{max. LP temporär um Rang}.
	\end{itemize}
	
	\medskip
	\textbf{Untote (skalieren mit Rang):}
	\begin{center}
		\begin{tabular}{ccccc}
			\toprule
			\textbf{Rang} & \textbf{Anzahl} & \textbf{LP} & \textbf{Schaden} & \textbf{Talente} \\
			\midrule
			1 & 1 & 10 & 1W4+2 & KAMPF 3, STÄRKE 1, GEWANDTHEIT 2 \\
			2 & 2 & 13 & 1W4+Rang & KAMPF 3, STÄRKE 2, GEWANDTHEIT 2 \\
			3 & 2 & 16 & 1W4+Rang & KAMPF 3, STÄRKE 3, GEWANDTHEIT 3 \\
			4 & 3 & 19 & 2W4+Rang & KAMPF 3, STÄRKE 4, GEWANDTHEIT 3 \\
			5 & 3 & 22 & 2W4+Rang & KAMPF 3, STÄRKE 5, GEWANDTHEIT 3 \\
			\bottomrule
		\end{tabular}
	\end{center}
	
	\textbf{Immunitäten:}
	Untote sind immun gegen \termref[Erosion]{erosion}, \termref[Blutung]{blutung} und \termref[Vergiftung]{vergiftung}.
}
{
	In einem speziellen Gegenstand können besiegte Gegner gespeichert werden:
	\begin{itemize}
		\item \textbf{2} normale Gegner, oder
		\item \textbf{1} normaler Gegner + \textbf{1} Miniboss (SL-Zustimmung), oder
		\item \textbf{1} Boss (SL-Zustimmung).
	\end{itemize}
	Es können nur ehemals lebendige Gegner gespeichert werden.
}

\paragraph*{Dunkler Fluch}

\spellcard
{\term{Dunkler Fluch}{dunkler-fluch}}
{Bonusaktion}
{8 MP}
{15 m}
{1 Stunde pro Rang oder bis aufgelöst}
{
	\begin{itemize}
		\item Ziel erleidet \textbf{+50\% Schaden}.
	\end{itemize}
	
	\medskip
	\textbf{Nebenwirkung:}
	\begin{itemize}
		\item Du erleidest selbst \textbf{+30\% Schaden}.
	\end{itemize}
}
{
	Ziel erhält zusätzlich \termref[Erosion]{erosion} Rang 2.
}

\paragraph*{Verfluchte Klinge}

\spellcard
{\term{Verfluchte Klinge}{verfluchte-klinge}}
{Aktion}
{8 MP}
{Selbst (eigene Waffe)}
{1 Runde pro Rang}
{
	\begin{itemize}
		\item Zielwaffe verursacht zusätzlich \textbf{+1W8 Schaden}.
	\end{itemize}
	
	\medskip
	\textbf{Nebenwirkung:}
	\begin{itemize}
		\item Bei Misslingen des Zaubers erleidet der Magier \textbf{1W8 Schaden}.
	\end{itemize}
}
{
	\textbf{Blutzehrung:}
	Solange der Zauber aktiv ist: Gib \textbf{1 Marke} aus, um mit jedem Treffer der Waffe \textbf{1W8 LP} zu erhalten (Überheilung möglich).
}

\paragraph*{Schattenklinge}

\spellcard
{\term{Schattenklinge}{schattenklinge}}
{Aktion}
{6 MP}
{Selbst}
{2 Runden}
{
	\begin{itemize}
		\item Beim Wirken greifst du sofort an.
		\item Die Klinge bleibt bis zum Ende deiner nächsten Runde bestehen.
		\item In Runde 2 kann der Magier seine \textbf{Aktion} für einen weiteren Angriff verwenden
		(\textbf{ohne} zusätzliche MP-Kosten).
	\end{itemize}
	
	\medskip
	\textbf{Schaden nach Rang:}
	\begin{center}
		\begin{tabular}{cc}
			\toprule
			\textbf{Rang} & \textbf{Schaden} \\
			\midrule
			1 & 1W10 \\
			2 & 1W10 + 2 \\
			3 & 1W10 + 4 \\
			4 & 1W10 + 6 \\
			5 & 1W10 + 8 \\
			6 & 1W10 + 10 \\
			\bottomrule
		\end{tabular}
	\end{center}
	
	\textbf{Nebenwirkungen:}
	\begin{itemize}
		\item Der Magier verliert bis zum Ende des Kampfes \textbf{Rang seiner Max. MP}.
		\item Bei Fehlschlag: Alle Verbündeten in \textbf{5 m} erleiden \textbf{2 + Rang Schaden},
		wenn ihre \textbf{STÄRKE-Probe} misslingt.
	\end{itemize}
}
{
	Gib \textbf{1 Marke} aus, um den Schaden zu \textbf{Direktschaden} zu machen.
}

\paragraph*{Schattenhaft}

\spellcard
{\term{Schattenhaft}{schattenhaft}}
{Bonusaktion}
{$3 + \text{Rang}$ MP}
{Selbst + (Rang--1) Verbündete in 5 m}
{1 Minute pro Rang}
{
	\begin{itemize}
		\item Du und bis zu \textbf{Rang--1} Verbündete werden zu Schattenwesen:
		\textbf{unsichtbar} und \textbf{beweglich}.
		\item Zum Zurückkehren wird eine \textbf{weitere Bonusaktion} benötigt.
		\item \textbf{Keine Angriffe} aus dem Schatten heraus möglich.
	\end{itemize}
	
	\medskip
	\textbf{Nebenwirkung:}
	\begin{itemize}
		\item Nicht-magiebegabte Mitreisende erleiden bei Rückkehr \textbf{1 + Rang Schaden}.
	\end{itemize}
}
{
	Bonusaktionen können als Angriff aus den Schatten heraus genutzt werden.
}


\spellcard
{\term{Seelenentzug}{seelenentzug}}
{Aktion}
{$4 + \text{Rang}$ MP}
{10 m}
{1 Runde pro Rang}
{
	\begin{itemize}
		\item Ziel erleidet pro Runde \textbf{Rang Schaden}.
		\item Du erhältst pro Runde \textbf{Rang LP} und \textbf{Rang MP} zurück.
		\item LP/MP dürfen \textbf{über das Maximum hinaus} wachsen.
	\end{itemize}
}
{
	Überschüssige LP können in \textbf{dauerhaftes} \termref[Magisches Schild (SP)]{magisches-schild} umgewandelt werden.
}

\subsection{Zeitmagie}

\subsubsection*{Rolle}
Chronomant / Manipulator des Zeitflusses

\subsubsection*{Kurzprofil}
Die \term{Zeitmagie}{klasse-zeitmagie} manipuliert den Ablauf der Realität selbst.
Zeitmagier verändern Aktionen, Runden, Initiativen und Zustände, statt rohen Schaden
zu verursachen. Ihre Magie ist mächtig, selten – und extrem riskant.

\subsubsection*{Stärken}
\begin{itemize}
	\item Kontrolle über Aktionen und Rundenverlauf
	\item Starke defensive und unterstützende Zauber
	\item Zugriff auf einzigartige Effekte (Rückspulen, Stasis)
\end{itemize}

\subsubsection*{Kampfstil}
Taktisch und vorausschauend – beeinflusst Initiativen, Aktionen,
Regeneration und gegnerische Bewegungen.

\subsubsection*{Magische Basisausrüstung (Empfehlung)}
\termref[Der Magischer Unterstützer]{ausruestung-magischer-unterstuetzer}

% =========================
% Klassenzauber
% =========================
\subsubsection*{Klassenzauber}

\paragraph*{Temporale Klinge}

\spellcard
{\term{Temporale Klinge}{temporale-klinge}}
{Aktion}
{6 MP}
{Nahkampf}
{2 Runden}
{
	\begin{itemize}
		\item \textbf{Sofort:} Verursacht \textbf{Rang + 2W4 Schaden}.
		\item \textbf{Nach 2 Runden:} Verursacht zusätzlich \textbf{Rang Schaden}
		durch Zeitenergie-Explosion.
	\end{itemize}
	
	\medskip
	\textbf{Nebenwirkung:}
	\begin{itemize}
		\item \textbf{Desorientierung:} \textbf{–1 Malus} auf die nächste Probe des Magiers
	\end{itemize}
}
{
	\textbf{Zeitschnitt:}
	Gib \textbf{1 Marke} aus:
	\begin{itemize}
		\item Schaden wird zu \termref[Durchschlagsschaden]{durchschlagsschaden}.
		\item Ziel verliert seine \textbf{nächste Bonusaktion}.
	\end{itemize}
}

\paragraph*{Zeitschub}

\spellcard
{\term{Zeitschub}{zeitschub}}
{Bonusaktion}
{$8 - \text{Rang}$ MP}
{10 m}
{Diese oder nächste Runde}
{
	\begin{itemize}
		\item Ziel erhält \textbf{1 zusätzliche Aktion}
		(in dieser oder der nächsten Runde).
	\end{itemize}
	
	\medskip
	\textbf{Nebenwirkung:}
	\begin{itemize}
		\item Ziel erhält \termref[Erschöpfung]{erschoepfung}.
	\end{itemize}
}
{
	Gib \textbf{1 Marke} aus:
	\begin{itemize}
		\item Ziel erhält \textbf{sofort 2 zusätzliche Bonusaktionen}.
	\end{itemize}
}

\paragraph*{Altersbrand}

\spellcard
{\term{Altersbrand}{altersbrand}}
{Aktion}
{10 MP}
{15 m}
{Sofort}
{
	\begin{itemize}
		\item \textbf{Objekte/Rüstungen:} Verlieren \textbf{80\%} ihrer Werte.
		\item \textbf{Magisches Schild (SP):} Verliert \textbf{100\%} der Werte.
		\item \textbf{Organische Wesen:}
		\begin{itemize}
			\item Erleiden \textbf{2 × Rang + Wurfzahl} \termref[Verwesungsschaden]{verwesungsschaden}
			\textbf{Rang}.
		\end{itemize}
	\end{itemize}
	
	\medskip
	\textbf{Nebenwirkung:}
	\begin{itemize}
		\item der Magier verliert \textbf{1W6 LP}.
		\item Mehrfaches Wirken im selben Kampf:
		\textbf{Nebenwirkung multipliziert sich}
		(z.\,B. 3.\ Einsatz = 3W6 LP).
	\end{itemize}
}
{
	Alle aktuell auf dem Ziel liegenden \termref[Statuseffekte]{statuseffekte}
	werden sofort abgehandelt.
}

\paragraph*{Zeit zurückdrehen}

\spellcard
{\term{Zeit zurückdrehen}{zeit-zurueckdrehen}}
{Bonusaktion}
{8 MP}
{10 m}
{Sofort}
{
	\begin{itemize}
		\item Dreht die Zeit für \textbf{1 Ziel} um bis zu
		\textbf{2 × Rang Sekunden} zurück.
		\item Geschehene Ereignisse in diesem Zeitraum werden für das Ziel rückgängig gemacht (z.B. Verlust/Erhalt von LP, MP, AP, Aktionen, Bonusaktionen).
	\end{itemize}
	
	\medskip
	\textbf{Hinweise:}
	\begin{itemize}
		\item \textbf{Rang 2:} Rücknahme einer Bonusaktion möglich.
		\item \textbf{Rang 3:} Rücknahme einer vollen Aktion möglich.
	\end{itemize}
}
{
	Gib \textbf{2 Marken} aus, um die Zeit für
	\textbf{Rang Ziele} zurückzudrehen.
}

\paragraph*{Tempoblase}

\spellcard
{\term{Tempoblase}{tempoblase}}
{Aktion}
{4 MP}
{Selbst (Radius Rang m)}
{1 Runde}
{
	\begin{itemize}
		\item \textbf{Projektilangriffe:} \textbf{50\%} Chance zu scheitern, wenn sie die Blase durchqueren.
		\item \textbf{Nahkampfangriffe:} \textbf{–2 Malus} für alle Ziele in der Blase.
		\item \textbf{Zauber:} Nur \textbf{50\% effektiv} (Schaden, Heilung, etc.).
	\end{itemize}
	
	\medskip
	\textbf{Nebenwirkung:}
	\begin{itemize}
		\item Ziele, die die Blase aktiv verlassen, erhalten \termref[Verwirrung]{verwirrung}.
	\end{itemize}
}
{
	\begin{itemize}
		\item Verbündete erhalten \textbf{2 zusätzliche Bonusaktionen}.
		\item Gegner verlieren 1 ihrer \textbf{Aktionen}.
	\end{itemize}
}

\paragraph*{Stasisfeld}

\spellcard
{\term{Stasisfeld}{stasisfeld}}
{Aktion}
{12 MP}
{10 m}
{Rang Runden oder nach Abbruch des Zaubers}
{
	\begin{itemize}
		\item Friert eine \textbf{1×1 m Zone} vollständig ein.
		\item Beim Eintritt:
		\begin{itemize}
			\item \textbf{Rang} \termref[Durchschlagsschaden]{durchschlagsschaden}.
			\item Ab \textbf{Wurfzahl 8}: zusätzlich \textbf{1W6 Schaden}.
		\end{itemize}
		\item Ziel kann keine Aktionen ausführen und sich nicht bewegen.
		\item Ziel ist \textbf{unverwundbar}.
		\item Status- und Zusatzeffekte wirken nicht,
		werden aber gespeichert und nach Ende von Stasisfeld normal abgehandelt.
	\end{itemize}
	
	\medskip
	\textbf{Nebenwirkung:}
	\begin{itemize}
		\item Solange der Zauber aktiv ist, kann der Magier keine Verteidigungsprobe ablegen
		(\textbf{automatischer Misserfolg}, aber der Angreifer muss trotzdem einen erfolgreichen Angriffswurf haben).
	\end{itemize}
}
{
	\begin{itemize}
		\item Pro Runde kann \textbf{1 Zeitfenster} geöffnet werden,
		um dem Ziel Schaden zuzufügen.
		\item Schadenseffekte über Zeit wirken nicht,
		können aber zugefügt werden.
	\end{itemize}
}
